% Options for packages loaded elsewhere
\PassOptionsToPackage{unicode}{hyperref}
\PassOptionsToPackage{hyphens}{url}
\PassOptionsToPackage{dvipsnames,svgnames,x11names}{xcolor}
\documentclass[
  brazilian,
  11pt,
  a4paper,
]{article}
\usepackage{xcolor}
\usepackage[margin=20mm]{geometry}
\usepackage{amsmath,amssymb}
\setcounter{secnumdepth}{-\maxdimen} % remove section numbering
\usepackage{iftex}
\ifPDFTeX
  \usepackage[T1]{fontenc}
  \usepackage[utf8]{inputenc}
  \usepackage{textcomp} % provide euro and other symbols
\else % if luatex or xetex
  \usepackage{unicode-math} % this also loads fontspec
  \defaultfontfeatures{Scale=MatchLowercase}
  \defaultfontfeatures[\rmfamily]{Ligatures=TeX,Scale=1}
\fi
\usepackage{lmodern}
\ifPDFTeX\else
  % xetex/luatex font selection
  \setmainfont[]{Candara}
\fi
% Use upquote if available, for straight quotes in verbatim environments
\IfFileExists{upquote.sty}{\usepackage{upquote}}{}
\IfFileExists{microtype.sty}{% use microtype if available
  \usepackage[]{microtype}
  \UseMicrotypeSet[protrusion]{basicmath} % disable protrusion for tt fonts
}{}
\makeatletter
\@ifundefined{KOMAClassName}{% if non-KOMA class
  \IfFileExists{parskip.sty}{%
    \usepackage{parskip}
  }{% else
    \setlength{\parindent}{0pt}
    \setlength{\parskip}{6pt plus 2pt minus 1pt}}
}{% if KOMA class
  \KOMAoptions{parskip=half}}
\makeatother
\usepackage{longtable,booktabs,array}
\newcounter{none} % for unnumbered tables
\usepackage{calc} % for calculating minipage widths
% Correct order of tables after \paragraph or \subparagraph
\usepackage{etoolbox}
\makeatletter
\patchcmd\longtable{\par}{\if@noskipsec\mbox{}\fi\par}{}{}
\makeatother
% Allow footnotes in longtable head/foot
\IfFileExists{footnotehyper.sty}{\usepackage{footnotehyper}}{\usepackage{footnote}}
\makesavenoteenv{longtable}
\usepackage{graphicx}
\makeatletter
\newsavebox\pandoc@box
\newcommand*\pandocbounded[1]{% scales image to fit in text height/width
  \sbox\pandoc@box{#1}%
  \Gscale@div\@tempa{\textheight}{\dimexpr\ht\pandoc@box+\dp\pandoc@box\relax}%
  \Gscale@div\@tempb{\linewidth}{\wd\pandoc@box}%
  \ifdim\@tempb\p@<\@tempa\p@\let\@tempa\@tempb\fi% select the smaller of both
  \ifdim\@tempa\p@<\p@\scalebox{\@tempa}{\usebox\pandoc@box}%
  \else\usebox{\pandoc@box}%
  \fi%
}
% Set default figure placement to htbp
\def\fps@figure{htbp}
\makeatother
\ifLuaTeX
  \usepackage{luacolor}
  \usepackage[soul]{lua-ul}
\else
  \usepackage{soul}
\fi
\ifLuaTeX
\usepackage[bidi=basic,shorthands=off]{babel}
\else
\usepackage[bidi=default,shorthands=off]{babel}
\fi
\ifPDFTeX
\else
\babelfont{rm}[]{Candara}
\fi
\ifLuaTeX
  \usepackage{selnolig} % disable illegal ligatures
\fi
\setlength{\emergencystretch}{3em} % prevent overfull lines
\providecommand{\tightlist}{%
  \setlength{\itemsep}{0pt}\setlength{\parskip}{0pt}}
\usepackage{bookmark}
\IfFileExists{xurl.sty}{\usepackage{xurl}}{} % add URL line breaks if available
\urlstyle{same}
\hypersetup{
  pdflang={pt-br},
  colorlinks=true,
  linkcolor={Maroon},
  filecolor={Maroon},
  citecolor={Blue},
  urlcolor={Blue},
  pdfcreator={LaTeX via pandoc}}

\author{}
\date{\vspace{-2.5em}}

\begin{document}

\part{Memorial Quantitativo}

\section{Identificação e
Formação}\label{identificauxe7uxe3o-e-formauxe7uxe3o}

\subsection{Dados Pessoais}\label{dados-pessoais}

\textbf{Nome completo:} André Luiz Barbosa Nunes da Cunha

\textbf{Nascimento:} 29/10/1981, Campo Grande/MS, Brasil

\textbf{Vínculo atual:} Professor Doutor (MS-3.2), Universidade de São
Paulo (USP), Escola de Engenharia de São Carlos (EESC), Departamento de
Engenharia de Transportes (STT)

\textbf{E-mail:}
\href{mailto:alcunha@usp.br}{\nolinkurl{alcunha@usp.br}}

\textbf{ORCID:}
\href{https://orcid.org/0000-0002-0520-0621}{0000-0002-0520-0621}

\textbf{Lattes:}
\href{http://lattes.cnpq.br/7996696632908127}{7996696632908127}

\textbf{Google Scholar:}
\href{https://scholar.google.com.br/citations?user=HI0CQJMAAAAJ}{André
Luiz Cunha}

\textbf{Idiomas:} Português (nativo), Inglês (avançado)

\subsection{Formação}\label{formauxe7uxe3o}

\begin{enumerate}
\def\labelenumi{\arabic{enumi}.}
\tightlist
\item
  \textbf{Doutorado em Ciências}, área de concentração Planejamento e
  Operação de Sistemas de Transportes, pela Universidade de São Paulo,
  Escola de Engenharia de São Carlos -- 2013.

  \begin{itemize}
  \tightlist
  \item
    \emph{Título}: Sistema automático para obtenção de parâmetros do
    tráfego veicular a partir de imagens de vídeo usando OpenCV
  \item
    \emph{Orientador}: Prof.~Titular José Reynaldo Anselmo Setti
  \item
    \emph{Bolsa}: Conselho Nacional de Desenvolvimento Científico e
    Tecnológico, CNPq
  \item
    \emph{DOI}:
    \href{https://doi.org/10.11606/T.18.2013.tde-19112013-165611}{10.11606/T.18.2013.tde-19112013-165611}
  \end{itemize}
\item
  \textbf{Mestrado em Engenharia de Transportes}, área de concentração
  Planejamento e Operação de Sistemas de Transportes, pela Universidade
  de São Paulo, Escola de Engenharia de São Carlos -- 2007.

  \begin{itemize}
  \tightlist
  \item
    \emph{Título}: Avaliação do impacto da medida de desempenho no
    equivalente veicular de caminhões
  \item
    \emph{Orientador}: Prof.~Titular José Reynaldo Anselmo Setti
  \item
    \emph{Bolsa}: Conselho Nacional de Desenvolvimento Científico e
    Tecnológico, CNPq
  \item
    \emph{DOI}:
    \href{https://doi.org/10.11606/D.18.2007.tde-27112007-094400}{10.11606/D.18.2007.tde-27112007-094400}.
  \end{itemize}
\item
  \textbf{Graduação em Engenharia Civil} pela Universidade Federal de
  Mato Grosso do Sul (UFMS), Campo Grande-MS, Brasil -- 2003.
\end{enumerate}

\subsection{Licenças}\label{licenuxe7as}

\begin{enumerate}
\def\labelenumi{\arabic{enumi}.}
\tightlist
\item
  Licença luto -- 2 a 9 de outubro de 2022.
\item
  Licença-paternidade -- 16 a 20 de maio de 2020.
\end{enumerate}

\subsection{Mobilidade Internacional}\label{mobilidade-internacional}

\begin{enumerate}
\def\labelenumi{\arabic{enumi}.}
\item
  \textbf{University of Melbourne (UoM)}, Austrália -- 2020 Pesquisador
  Visitante Junior, bolsa CAPES-Print.
\item
  \textbf{University of Zagreb (UNIZG)}, Croácia -- 2018 e 2022
  \emph{Professor Visitante}, programa ERASMUS+.
\item
  \textbf{Universidade do Minho (UMinho)}, Portugal -- 2017
  \emph{Professor Visitante}, programa de mobilidade CAPES-FCT.
\item
  \textbf{Technical University of Munich (TUM)}, Alemanha -- 2016 e 2017
  \emph{Pesquisador Visitante}, programa de mobilidade FAPESP-BAYLAT.
\end{enumerate}

\subsection{Prêmios}\label{pruxeamios}

\begin{enumerate}
\def\labelenumi{\arabic{enumi}.}
\item
  \textbf{Prêmio ANPET de Produção Científica}, Agência Nacional de
  Pesquisa e Ensino em Transportes (ANPET) -- 2023.
\item
  \textbf{Certificado de Excelência}, melhor docente do Departamento de
  Engenharia de Transportes (USP-EESC-STT), SACivil - Secretaria
  Acadêmica da Engenharia Civil -- 2017.
\item
  \textbf{Certificado de Excelência}, melhor docente do Departamento de
  Engenharia de Transportes (USP-EESC-STT), SACivil - Secretaria
  Acadêmica da Engenharia Civil -- 2016.
\item
  \textbf{Prêmio Salão de Inovação ABCR}, 9º Congresso Brasileiro de
  Rodovias e Concessões (CBR\&C), 5º Salão de Inovação da Associação
  Brasileira de Concessionária de Rodovias (ABCR) -- 2015.
\end{enumerate}

\subsection{Formação complementar}\label{formauxe7uxe3o-complementar}

\begin{enumerate}
\def\labelenumi{\arabic{enumi}.}
\item
  \textbf{EMI for Professors}. Curso de capacitação docente (Carga
  horária: 80h) em \emph{English as a Medium of Instruction} (EMI).
  Universidade de São Paulo, AUCANI, São Paulo, SP, Brasil. --
  18/08/2025 a 12/12/2025
\item
  \textbf{Microssimulação de Ciclovias no PTV Vissim}. Curso de curta
  duração (Carga horária: 2h), Congresso ANPET (Associação Nacional de
  Pesquisa e Ensino em Transporte), Florianópolis, SC, Brasil --
  06/11/2024.
\item
  \textbf{Um curso prático sobre assessibilidade urbana em R}. Curso de
  curta duração (Carga horária: 2h), Congresso ANPET (Associação
  Nacional de Pesquisa e Ensino em Transporte), Florianópolis, SC,
  Brasil -- 06/11/2024.
\item
  \textbf{Exploração e processamento conjunto de dados públicos para
  pesquisa e prática de tráfego veicular}. Curso de curta duração (Carga
  horária: 1,5h), Congresso ANPET (Associação Nacional de Pesquisa e
  Ensino em Transporte), Florianópolis, SC, Brasil -- 05/11/2024.
\item
  \textbf{Ensino Superior: compartilhando conhecimentos entre docentes
  USP}. Pró-reitoria de Pesquisa da USP (Carga horária: 15h), Palestras
  online, Brasil -- 3 a 12 de março de 2021.
\item
  \textbf{\emph{Design thinking}, laboratórios de simulação e
  gameficação como estratégias para o ensino - novas soluções para
  velhos problemas}. Desenvolvimento e Capacitação Docente da USP (Carga
  horária: 4h), Palestras online, Brasil -- 07/10/2019.
\item
  \textbf{Metodologias ativas na inovação do ensino superior}.
  Desenvolvimento e Capacitação Docente da USP (Carga horária: 4h),
  Palestras online, Brasil -- 03/10/2019.
\item
  \textbf{\emph{São Paulo School of Advanced Science on Learning from
  Data} (SPSAS)}. Curso de curta duração (Carga horária: 80h),
  Universidade de São Paulo (USP), São Paulo, SP, Brasil -- 29 de julho
  a 9 de agosto de 2019.
\item
  \textbf{Disciplinar, Interdisciplinar e o Ensino por Projetos}.
  Seminário (Carga horária: 1,5h), Instituto de Química de São Carlos
  (USP-IQSC), São Carlos, SP, Brasil -- 12/06/2019.
\item
  \textbf{Curso de Python 3 - Fundamentos}. Curso de curta duração
  online (Carga horária: 40h), Curso em Vídeo, CEV, Brasil --
  17/10/2017.
\item
  \textbf{Como ensinar hoje: modelos blended, metodologias ativas e
  tecnologias digitais}. Ciclo de Atividades de Desenvolvimento Docente
  da USP (Carga horária: 2h), Ribeirão Preto, SP, Brasil -- 15/09/2017.
  {[}\textbf{{[}@Fig:USP-mativas{]}}{]}
\item
  \textbf{Compreendendo Deep Learning}. Curso de curta duração (Carga
  horária: 8h), Neural Mind Technologies (NMT), Campinas, SP, Brasil --
  01/09/2017.
\item
  \textbf{Escola Avançada em Big Data Analysis}. Cursos de curta duração
  (Carga horária: 24h), Instituto de Ciências Matemáticas e de
  Computação, ICMC, Brasil -- 3 a 7 de julho de 2017.

  \begin{itemize}
  \tightlist
  \item
    ``\emph{Deep Learning}'' (Carga horária: 8h)
  \item
    ``\emph{Visualização de dados}'' (Carga horária: 8h)
  \item
    ``\emph{Aprendizado de Máquinas}'' (Carga horária: 8h).
  \end{itemize}
\item
  \textbf{II Workshop em Ciências de Dados}. Curso de curta duração
  (Carga horária: 4h), Universidade de São Paulo, USP, São Paulo, Brasil
  -- 11/05/2017.
\item
  \textbf{Uso da Plataforma Turnitin}. Curso de curta duração (Carga
  horária: 3h), Escola de Engenharia de São Carlos, USP-EESC, Brasil --
  09/05/2017.
\item
  \textbf{Minicurso de Arduino}. Extensão universitária (Carga horária:
  5h), Escola de Engenharia de São Carlos, USP-EESC, Brasil --
  01/10/2016.
\item
  \textbf{DS101X: Statistical Thinking for Data Science and Analytics}.
  Curso de curta duração online (Carga horária: 12h), Columbia
  University, New York, Estados Unidos -- 22/01/2016.
\item
  \textbf{Using Python to Access Web Data}. Curso de curta duração
  online (Carga horária: 20h), University of Michigan, UMICH, Ann Arbor,
  Estados Unidos -- 2015.
\item
  \textbf{Getting Started with Python}. Curso de curta duração online
  (Carga horária: 20h), University of Michigan, UMICH, Ann Arbor,
  Estados Unidos -- 2015.
\item
  \textbf{Estruturas de Dados - Python}. Curso de curta duração online
  (Carga horária: 12h), University of Michigan, UMICH, Ann Arbor,
  Estados Unidos -- 2015.
\item
  \textbf{Utilização de Bancos de Dados com Python}. Curso de curta
  duração online (Carga horária: 12h), University of Michigan, UMICH,
  Ann Arbor, Estados Unidos -- 2015.
\item
  \textbf{EX101x Data Analysis: Take it to the MAX()}. Curso de curta
  duração online (Carga horária: 40h), Delft University of Technology,
  TU DELFT, Delft, Holanda -- 02/12/2015.
\item
  \textbf{Elaboração de ebooks para ensino digital, portátil e
  flexível}. Ciclo de Atividades de Desenvolvimento Docente da USP
  (Carga horária: 2h), Ribeirão Preto, SP, Brasil -- 12/11/2015.
\item
  \textbf{DAT203x Data Science \& Machine Learning Essential}. Curso de
  curta duração online (Carga horária: 20h), Microsoft Corporation,
  Washington, D. C., Estados Unidos -- 01/11/2015.
\item
  \textbf{Criação de vídeos-tutoriais de uma forma prática}. Ciclo de
  Atividades de Desenvolvimento Docente da USP (Carga horária: 2h),
  Ribeirão Preto, SP, Brasil -- 15/10/2015.
\item
  \textbf{Estética e dinâmica para a elaboração de aulas eletrônicas}.
  Ciclo de Atividades de Desenvolvimento Docente da USP (Carga horária:
  2h), Ribeirão Preto, SP, Brasil -- 24/09/2015.
\item
  \textbf{Da realidade virtual à fabricação digital: Desafios e
  Oportunidades em sala de aula}. Ciclo de Atividades de Desenvolvimento
  Docente da USP (Carga horária: 2h), Ribeirão Preto, SP, Brasil --
  10/09/2015.
\item
  \textbf{Curso de PHP}. Curso de curta duração online (Carga horária:
  40h), CursoemVideo.com, Brasil -- 03/05/2015.
\item
  \textbf{Processamento de Imagens usando Python e NumPy}. Extensão
  universitária (Carga horária: 20h), Universidade Estadual de Campinas,
  UNICAMP, Campinas, Brasil -- 12 de março a 11 de abril de 2015.
\item
  \textbf{Machine Learning Summer School} (MLSS). Curso de curta duração
  (Carga horária: 30h), Sydney, New South Wales, Austrália -- 16 a 25 de
  fevereiro de 2015.
\end{enumerate}

\subsection{Experiência Profissional}\label{experiuxeancia-profissional}

\begin{enumerate}
\def\labelenumi{\arabic{enumi}.}
\item
  \textbf{Prestação de serviço de consultoria para a empresa Canhedo
  Beppu Engenheiros Associados Ltda} para a análise do comportamento da
  frota de caminhões ao longo da nova pista descendente da Via Dutra
  (BR-116), no trecho da Serra das Araras, no Estado do Rio de Janeiro
  sob jurisdição da Motiva CCR RioSP Nova Dutra -- agosto a dezembro de
  2025.
\item
  \textbf{Prestação de serviço de consultoria para a empresa Canhedo
  Beppu Engenheiros Associados Ltda} para a análise da localização de
  áreas de escape para caminhões sem freios no projeto da nova pista
  descendente da Via Dutra (BR-116), no trecho da Serra das Araras, no
  Estado do Rio de Janeiro sob jurisdição da CCR Nova Dutra -- julho a
  novembro de 2023.
\item
  \textbf{Professor Visitante (virtual) na University of Zagreb},
  Zagreb, Croácia. Programa ERASMUS+: Mobilidade Docente para Ensino
  Virtual (online) - Acordo de Mobilidade. Palestras online para cursos
  de graduação e pós-graduação na University of Zagreb, Faculty of
  Transport and Traffic Sciences -- 11 a 15 de abril de 2022.
\item
  \textbf{Anfitrião na USP-EESC da mobilidade virtual do Prof.~Edouard
  Ivanjko}, Zagreb, Croácia. Programa ERASMUS+: Mobilidade Docente para
  Ensino Virtual (online) - Acordo de Mobilidade. Palestras online para
  cursos de graduação e pós-graduação na University of Zagreb, Faculty
  of Transport and Traffic Sciences (Carga horária: 8h) -- 11 a 15 de
  abril de 2022.
\item
  \textbf{Professor Visitante na University of Melbourne}, Melbourne
  School of Engineering, Department of Infrastructure Engineering, sob
  supervisão dos professores Majid Sarvi e Patricia Sauri Lavieri. O
  projeto ``\emph{Uma abordagem baseada em grafos para análise
  espaço-temporal de redes viárias}'' foi contemplado com uma bolsa de
  um ano pela CAPES (processo CAPES-Print -- 88887.371506/2019-00) pelo
  Programa Print - Professor Visitante Junior -- janeiro a dezembro de
  2020.
\item
  \textbf{Assessoria junto a empresa Autopista Litoral Sul} (APLS S.A.),
  grupo ARTERIS, em que foi realizada a supervisão de ensaios de
  homologação na área de escape para caminhões no km 667 da BR-376/PR --
  dezembro de 2019.
\item
  \textbf{Professor Anfitrião da USP-EESC} ao receber a visita do prof.
  Eduard Ivanjko, Universidade de Zagreb, Zagreb, Croácia. Programa
  ERASMUS+: Acordo de Mobilidade no Ensino Superior (UNIZG / USP-EESC).
  Mobilidade docente para compartilhar experiências de ensino e pesquisa
  em Sistemas Inteligentes de Transportes (ITS). O professor visitante
  realizou palestras na USP-EESC e na USP-Politécnica, em São Paulo
  (Carga horária: 8h) -- 2 a 8 de abril de 2019.
\item
  \textbf{Ministrou a disciplina ``\emph{Tecnologia dos Transportes
  II}'' no Curso de Capacitação em Infraestrutura de Transportes}, pelo
  convênio entre a USP e o Ministério de Obras Públicas e Comunicações
  do Paraguai (MOPC). (Carga horária: 15h) -- 1 de agosto a 30 de
  novembro de 2018.
\item
  \textbf{Mediador e palestrante no evento ``\emph{Arena New Mobility -
  Salão de Automóveis 2018}''} realizado em São Paulo, SP, Brasil. O
  docente foi mediador em duas mesas (Carga horária: 4,5h) -- 15 e 16 de
  novembro de 2018.
\item
  \textbf{Professor Visitante na Universidade de Zagreb}, Zagreb,
  Croácia. Programa ERASMUS+: Acordo de Mobilidade no Ensino Superior
  (UNIZG / USP-EESC). Mobilidade docente para compartilhar experiências
  de ensino e pesquisa em Sistemas Inteligentes de Transportes (ITS).
  Cursos ministrados para os níveis de graduação e pós-graduação na
  University of Zagreb, Faculty of Transport and Traffic Sciences (Carga
  horária: 13h) -- 4 a 14 de junho de 2018.
\item
  \textbf{Ministrou a disciplina ``\emph{Sistemas Inteligentes de
  Transportes Rodoviários e Urbanos}'' no Curso de Capacitação em
  Infraestrutura de Transportes}, pelo convênio entre a USP e o
  Ministério de Obras Públicas e Comunicações do Paraguai (MOPC). (Carga
  horária: 32h) -- 1 de março a 1 de julho de 2018.
\item
  \textbf{Professor Visitante no \emph{TUM-USP Workshop on Sustainable
  Mobility}} realizado na USP Escola Politécnica, São Paulo, SP, Brasil.
  O workshop envolveu pesquisadores alemães e brasileiros selecionados
  pela Chamada do edital BAYLAT/FAPESP -- 18 a 21 de setembro de 2017.
\item
  \textbf{Professor Visitante na Universidade do Minho}, Guimarães,
  Portugal. Esta missão internacional fez parte do projeto aprovado
  CAPES-FCT 39/2014 (realizado entre 2016 e 2018). Durante a missão, foi
  realizada a visita às instituições Universidade do Minho e
  Universidade do Porto para discussão de projetos colaborativos no tema
  de caracterização do comportamento de pedestres -- 17 a 28 de julho de
  2017.
\item
  \textbf{Professor Visitante no \emph{TUM-USP Workshop on Sustainable
  Mobility}} realizado no Institute of Automotive Technology (FTM),
  Technical University of Munich (TUM), em Munique, Bavaria, Alemanha. O
  workshop envolveu pesquisadores alemães e brasileiros selecionados
  pela Chamada do edital BAYLAT/FAPESP -- 28 de novembro a 2 de dezembro
  de 2016.
\item
  \textbf{Professor Doutor (MS-3.1) no Departamento de Engenharia de
  Transportes (STT)}, Escola de Engenharia de São Carlos (USP-EESC), São
  Carlos, SP, Brasil. Concurso Público pelo Edital ATAc-56/2013 (Carga
  Horária: Regime de Dedicação Integral à Docência e à Pesquisa - RDIDP)
  -- desde 1 de julho de 2014 até o presente momento.
\item
  \textbf{Especialista em Laboratório (Superior 1A) no Departamento de
  Engenharia de Transportes (STT)}, Escola de Engenharia de São Carlos
  (USP-EESC), São Carlos, SP, Brasil. Concurso Público pelo Edital
  EESC/USP 11/2012 (Carga horária: 40h/semana). \emph{Atividades
  realizadas}: (1) Desenvolver pesquisas científicas em projetos
  liderados por docentes, com foco didático-científico e de extensão;
  (2) Coletar, manipular e publicar dados de pesquisa para disseminação
  do conhecimento; (3) Auxiliar na orientação de alunos de iniciação
  científica e pós-graduação; (4) Apoiar na organização de atividades em
  laboratórios; (5) Dar suporte a docentes em pesquisa e extensão,
  excluindo atividades didáticas -- 14 de fevereiro de 2013 a 30 de
  junho de 2014.
\item
  \textbf{Consultor Externo em Engenharia de Transportes no Técnicos em
  Transportes Ltda (TECTRAN)}, Belo Horizonte, MG, Brasil.
  \emph{Atividades realizadas}: Criação de banco de dados do EPELT
  (Escritório de Planejamento de Logística de Transportes) da Secretaria
  do Estado de Minas Gerais -- 2 de abril de 2012 a 1 de dezembro de
  2012.
\item
  \textbf{Prestação de serviço de Engenharia Civil como Autônomo no
  Instituto de Ciências Matemáticas e da Computação (ICMC)}, da USP. São
  Carlos, SP, Brasil. \emph{Atividades realizadas}: Digitalização da
  infraestrutura em AutoCAD, manutenção de edifícios e gestão das obras
  em andamento no ICMC -- 1 de março de 2012 a 1 de abril de 2012.
\item
  \textbf{Professor Substituto na Universidade Estadual Paulista (UNESP)
  ``Júlio de Mesquita Filho'', Faculdade de Engenharia de Bauru (FEB)}.
  Bauru, SP, Brasil. \emph{Atividades realizadas}: Ministrei duas
  disciplinas do curso da graduação em Engenharia Civil: ``\emph{Técnica
  e Economia dos Transportes}'' no 1º semestre; e ``\emph{Engenharia de
  Tráfego}'' no 2º semestre. (Carga horária: 12h/semana) -- 2 de março
  de 2010 a 20 de dezembro de 2010.
\item
  \textbf{Estágio Supervisionado em Docência, bolsista do Programa de
  Aperfeiçoamento do Ensino (PAE)}, da Universidade de São Paulo, junto
  à disciplina de graduação \emph{STT0408 -- Introdução à Engenharia de
  Transportes}, do Departamento de Transportes da Escola de Engenharia
  de São Carlos, São Carlos, SP, Brasil. (Carga horária: 6h/semana) --
  fevereiro a junho de 2009.
\item
  \textbf{Prestação de serviço de editoração de texto da Revista
  TRANSPORTES}, Rio de Janeiro, RJ, Brasil -- junho de 2006 a novembro
  de 2011.
\item
  \textbf{Estágio Supervisionado em Docência, bolsista do Programa de
  Aperfeiçoamento do Ensino (PAE)}, da Universidade de São Paulo, junto
  à disciplina de graduação \emph{STT0408 -- Introdução à Engenharia de
  Transportes}, do Departamento de Transportes da Escola de Engenharia
  de São Carlos, São Carlos, SP, Brasil. Carga horária de 6h/semana --
  fevereiro a junho de 2006.
\item
  \textbf{Estagiário em Engenharia Civil no Comércio, Construção e
  Engenharia Ltda (COCENG)}. Campo Grande, MS, Brasil. \emph{Atividades
  realizadas}: Desenhos em AutoCAD de projetos residenciais, orçamento
  de obra e fiscalização de obras -- 1 de março de 2003 a 2 de fevereiro
  de 2004.
\item
  \textbf{Estagiário em Engenharia Civil no Desenvolvimento de Projetos
  de Engenharia Ltda (DPE)}. Campo Grande, MS, Brasil. \emph{Atividades
  realizadas}: Desenhos em AutoCAD de projetos geométrico de rodovias e
  de sinalização horizontal e vertical -- 2 de dezembro de 2002 a 28 de
  fevereiro de 2003.
\end{enumerate}

\section{Trajetória Docente e
Pedagógica}\label{trajetuxf3ria-docente-e-pedaguxf3gica}

\subsection{Disciplinas Ministradas}\label{disciplinas-ministradas}

Em síntese, as disciplinas ministradas podem ser organizadas nos
seguintes eixos:

\begin{enumerate}
\def\labelenumi{\arabic{enumi}.}
\tightlist
\item
  \textbf{Ciência de Dados e Engenharia Computacional}

  \begin{itemize}
  \tightlist
  \item
    \textbf{Análise Multivariada de Dados e Estatística Aplicada}
    (Pós-Graduação) Modernizou o currículo da pós-graduação para incluir
    técnicas de IA (Redes Neurais, Árvores de Decisão, Análise de
    Cluster, Algoritmos Genéticos) e programação em R. Orientou
    estudantes na publicação de pesquisas originais a partir dos
    conjuntos de dados trabalhados na disciplina.
  \item
    \textbf{Ferramentas Computacionais para Engenharia Civil}
    (Graduação) Concebeu e implementou nova disciplina eletiva para
    suprir a lacuna de competências digitais, programação em planilhas,
    resolução algorítmica de problemas, AutoCAD, Civil 3D e SIG.
  \end{itemize}
\item
  \textbf{Engenharia de Transportes e Simulação}

  \begin{itemize}
  \tightlist
  \item
    \textbf{Fundamentos de Engenharia de Transportes} (Graduação)
    Alcançou avaliações discentes excepcionais (4,5/5,0) mediante a
    adoção de estratégias de sala de aula invertida e aprendizagem
    baseada em projetos para turmas de mais de 50 estudantes.
  \item
    \textbf{Engenharia de Tráfego e Simulação} (Graduação/Pós-Graduação)
    Oferece formação avançada em teorias de simulação microscópica
    (TSIS-CORSIM, VISSIM, AIMSUN e SUMO), preparando os estudantes com
    competências técnicas prontas para o mercado.
  \end{itemize}
\item
  \textbf{Logística e Infraestrutura}

  \begin{itemize}
  \tightlist
  \item
    \textbf{Logística e Gestão da Cadeia de Suprimentos} Reformulou
    integralmente a ementa para responder a desafios contemporâneos,
    integrando ferramentas de logística baseadas em IA, planejamento de
    rotas com SIG e boas práticas de Logística Verde.
  \item
    \textbf{Aeroportos e Hidrovias} Gerenciou disciplinas obrigatórias
    de grande porte (mais de 50 estudantes), equilibrando fundamentos
    teóricos com exercícios práticos de operação.
  \end{itemize}
\item
  \textbf{Projetos de Conclusão de Curso e Mentoria}

  \begin{itemize}
  \tightlist
  \item
    \textbf{Projetos de Conclusão de Curso} (Graduação) Supervisionou
    mais de 25 projetos, conduzindo os estudantes desde a definição do
    problema de pesquisa até a elaboração de relatórios técnicos de
    padrão profissional.
  \item
    \textbf{Supervisão de Estágios} (Graduação) Acompanha inserções
    profissionais de mais de 15 estudantes, assegurando a articulação
    entre os objetivos acadêmicos e as demandas do setor.
  \end{itemize}
\end{enumerate}

\subsubsection{Graduação}\label{graduauxe7uxe3o}

Os códigos e nomes das disciplinas oferecidas e ministradas pelo docente
foram:

\begin{enumerate}
\def\labelenumi{\arabic{enumi}.}
\tightlist
\item
  1800078 - Estágio Supervisionado em Engenharia Civil
\item
  1800093 - Trabalho de Conclusão de Curso I
\item
  1800094 - Projeto Final de Curso II
\item
  1800122 - Estágio Supervisionado
\item
  1800123 - Desenho
\item
  STT0403 - Aeroportos, Portos e Vias Navegáveis
\item
  STT0408 - Fundamentos de Engenharia de Transportes
\item
  STT0412 - Ferramentas computacionais aplicadas à Engenharia Civil
  (\emph{atualmente STT0630})
\item
  STT0602 - Engenharia de Tráfego
\item
  STT0610 - Logística e Transportes (\emph{cancelada e criada a
  STT0631})
\item
  STT0618 - Transporte Aéreo
\item
  STT0628 - Engenharia de Tráfego e Simulação de Tráfego Rodoviário
\item
  STT0630 - Ferramentas computacionais aplicadas à Engenharia Civil
\item
  STT0631 - Logística em obras civis (\emph{início do oferecimento 2º
  semestre de 2026})
\end{enumerate}

A Tabela @tab:grad resume a carga horária (h/semana) e o oferecimento a
cada semestre de 2014 a 2026.

\begin{longtable}[]{@{}
  >{\centering\arraybackslash}p{(\linewidth - 6\tabcolsep) * \real{0.1343}}
  >{\centering\arraybackslash}p{(\linewidth - 6\tabcolsep) * \real{0.1642}}
  >{\raggedright\arraybackslash}p{(\linewidth - 6\tabcolsep) * \real{0.5522}}
  >{\centering\arraybackslash}p{(\linewidth - 6\tabcolsep) * \real{0.1493}}@{}}
\caption{Carga horária de disciplinas na graduação por semestre
(2014-2026).}\label{tab:grad}\tabularnewline
\toprule\noalign{}
\begin{minipage}[b]{\linewidth}\centering
Sem.
\end{minipage} & \begin{minipage}[b]{\linewidth}\centering
Ano
\end{minipage} & \begin{minipage}[b]{\linewidth}\raggedright
Disciplinas
\end{minipage} & \begin{minipage}[b]{\linewidth}\centering
h/semana
\end{minipage} \\
\midrule\noalign{}
\endfirsthead
\toprule\noalign{}
\begin{minipage}[b]{\linewidth}\centering
Sem.
\end{minipage} & \begin{minipage}[b]{\linewidth}\centering
Ano
\end{minipage} & \begin{minipage}[b]{\linewidth}\raggedright
Disciplinas
\end{minipage} & \begin{minipage}[b]{\linewidth}\centering
h/semana
\end{minipage} \\
\midrule\noalign{}
\endhead
\bottomrule\noalign{}
\endlastfoot
2º & 2014 & STT0618 & 1:40 \\
1º & 2015 & STT0403, STT0408 & 7:18 \\
2º & 2015 & STT0602 & 1:57 \\
1º & 2016 & 1800093, STT0403, STT0408 & 9:26 \\
2º & 2016 & STT0412, STT0628 & 5:57 \\
1º & 2017 & STT0403, STT0408 & 8:30 \\
2º & 2017 & STT0412 & 2:19 \\
1º & 2018 & STT0408 & 5:19 \\
2º & 2018 & STT0412, STT0628 & 6:02 \\
1º & 2019 & 1800078, 1800093, STT0403, STT0408 & 11:36 \\
2º & 2019 & 1800078, STT0412, STT0628 & 7:36 \\
--- & 2020 & \emph{Afastamento na Univ. of Melbourne} & --- \\
1º & 2021 & 1800078, STT0403, STT0408 & 7:38 \\
2º & 2021 & 1800078, STT0628, STT0630 & 9:04 \\
1º & 2022 & 1800078, STT0403, STT0408 & 9:01 \\
2º & 2022 & STT0628, STT0630 & 6:06 \\
1º & 2023 & 1800078, STT0403, STT0408 & 8:25 \\
2º & 2023 & 1800094, 1800122, STT0628, STT0630 & 9:15 \\
1º & 2024 & 1800094, STT0403, STT0408 & 12:52 \\
2º & 2024 & 1800122, STT0628, STT0630 & 7:36 \\
1º & 2025 & STT0408, STT0610 & 6:36 \\
2º & 2025 & 1800122, STT0628 & 5:24 \\
1º & 2026 & 1800123, STT0408 & 10:45 \\
\textbf{Total} & \textbf{22 sem.} & \textbf{52 oferecimentos · 13
disciplinas} & \textbf{160:22} \\
& & & \\
\end{longtable}

\subsubsection{Pós-graduação}\label{puxf3s-graduauxe7uxe3o}

As disciplinas oferecidas pelo docente na pós-graduação foram:

\begin{enumerate}
\def\labelenumi{\arabic{enumi}.}
\tightlist
\item
  STT5859 - Tecnologia dos Transportes (1º semestre)
\item
  STT5874 - Tópicos Avançados de Engenharia de Tráfego (2º semestre)
\item
  STT5898 - Estatística Aplicada à Engenharia de Transportes (1º
  semestre)
\item
  STT5900 - Análise de Dados Multivariados Aplicados à Engenharia de
  Transportes (2º semestre)
\item
  STT5905 - Pesquisa Bibliográfica para Planejamento e Operação de
  Sistemas de Transportes (2º semestre)
\item
  STT5909 - Laboratório de Análise de Dados com Software Livre R.
  (\emph{oferecimento cancelado, pois a linguagem R foi incorporada a
  disciplina STT5900})
\end{enumerate}

A Tabela @tab:posgrad resume a carga horário semanal em disciplinas da
pós-graduação desde o início do credenciamento do docente em 2015 até
2026.

\begin{longtable}[]{@{}
  >{\centering\arraybackslash}p{(\linewidth - 4\tabcolsep) * \real{0.1286}}
  >{\raggedright\arraybackslash}p{(\linewidth - 4\tabcolsep) * \real{0.7429}}
  >{\centering\arraybackslash}p{(\linewidth - 4\tabcolsep) * \real{0.1286}}@{}}
\caption{Carga horária de disciplinas na pós-graduação por ano
(2014-2026).}\label{tab:posgrad}\tabularnewline
\toprule\noalign{}
\begin{minipage}[b]{\linewidth}\centering
Ano
\end{minipage} & \begin{minipage}[b]{\linewidth}\raggedright
Disciplinas
\end{minipage} & \begin{minipage}[b]{\linewidth}\centering
h/semana
\end{minipage} \\
\midrule\noalign{}
\endfirsthead
\toprule\noalign{}
\begin{minipage}[b]{\linewidth}\centering
Ano
\end{minipage} & \begin{minipage}[b]{\linewidth}\raggedright
Disciplinas
\end{minipage} & \begin{minipage}[b]{\linewidth}\centering
h/semana
\end{minipage} \\
\midrule\noalign{}
\endhead
\bottomrule\noalign{}
\endlastfoot
2015 & STT5874, STT5898, STT5900 & 6:00 \\
2016 & STT5859, STT5874, STT5898, STT5900 & 7:20 \\
2017 & STT5859, STT5874, STT5898, STT5900, STT5905, STT5909 & 9:40 \\
2018 & STT5859, STT5874, STT5898, STT5900, STT5905 & 8:00 \\
2019 & STT5859, STT5874, STT5898, STT5900, STT5905 & 8:00 \\
2020 & \emph{Afastamento na Univ. of Melbourne} & --- \\
2021 & STT5859, STT5874, STT5898, STT5900, STT5905 & 7:24 \\
2022 & STT5859, STT5874, STT5898, STT5900 & 7:00 \\
2023 & STT5859, STT5874, STT5898, STT5900, STT5905 & 7:24 \\
2024 & STT5859, STT5874, STT5898, STT5900, STT5905 & 7:24 \\
2025 & STT5859, STT5874, STT5898, STT5900, STT5905 & 7:24 \\
2026 & STT5859, STT5898 & 3:00 \\
\textbf{Total} & \textbf{49 oferecimentos · 6 disciplinas} &
\textbf{78:36} \\
\end{longtable}

\subsection{Orientações
Acadêmicas}\label{orientauxe7uxf5es-acaduxeamicas}

O docente supervisionou mais de 60 pesquisadores em todos os níveis
acadêmicos:

\begin{itemize}
\tightlist
\item
  \textbf{Doutorado}: 6 no total (3 concluídos, 3 em andamento)
\item
  \textbf{Mestrado}: 16 no total (12 concluídos, 4 em andamento)
\item
  \textbf{Iniciação Científica}: 18 no total (16 concluídas, 2 em
  andamento)
\item
  \textbf{Trabalhos de Conclusão de Curso}: 26 concluídos.
\end{itemize}

\textbf{Principais Teses e Dissertações Orientadas}

\begin{description}
\tightlist
\item[IA e Visão Computacional]
Supervisionou aplicações inéditas de Deep Learning e Modelos de
Linguagem Multimodal (LMMs) para detecção automática de eixos de
caminhões e classificação de veículos.
\item[Resiliência de Infraestrutura]
Orientou pesquisas sobre aplicações de Teoria dos Grafos para mensuração
da vulnerabilidade de redes urbanas a eventos extremos de inundação.
\item[Mobilidade Sustentável]
Supervisiona pesquisa de doutorado em andamento sobre os impactos das
bicicletas elétricas na equidade espacial e na acessibilidade a empregos
em áreas metropolitanas.
\end{description}

\textbf{Inserção Profissional dos Egressos}

Os egressos do grupo de pesquisa ocupam posições estratégicas em
Agências Reguladoras Nacionais (como ANTT e ARTESP), outros em
consultorias líderes de Engenharia de Transportes, empresas de
transportes (LATAM) e alguns no setor financeiro.

\subsubsection{Orientações
Concluídas}\label{orientauxe7uxf5es-concluuxeddas}

\paragraph{Doutorado (3)}\label{doutorado-3}

\begin{enumerate}
\def\labelenumi{\arabic{enumi}.}
\item
  \href{https://scholar.google.com/citations?user=AcNWryQAAAAJ}{Leandro
  Arab Marcomini}. \textbf{Classificação de caminhões por eixos com
  aprendizagem profunda e um modelo multimodal de larga escala}. 2025.
  Tese (Engenharia de Transportes) - Universidade de São Paulo. Bolsa:
  CNPq\\
  DOI:
  \href{https://doi.org/10.11606/T.18.2025.tde-11022026-082130}{10.11606/T.18.2025.tde-11022026-082130}
\item
  \href{https://scholar.google.com/citations?user=BqdPZ2gAAAAJ}{Andre
  Borgato Morelli}. \textbf{Análise da vulnerabilidade em redes urbanas
  brasileiras: Explorando o impacto de alagamentos e sua relação com a
  morfologia e uso do solo}. 2025. Tese (Engenharia de Transportes) -
  Universidade de São Paulo. Bolsa: CAPES\\
  DOI: \href{}{ainda sem a referência no teses.usp.br}
\item
  \href{https://lattes.cnpq.br/1828432370828925}{Gabriel Jurado Martins
  de Oliveira}. \textbf{Toward a data-driven framework to analyse
  temporal and spatial distributions of crashes}. 2024. Tese (PhD in
  Infrastructure Engineering) - The University of Melbourne
  (co-orientador). Acesso em: \url{https://hdl.handle.net/11343/354471}
\end{enumerate}

\paragraph{Mestrado (12)}\label{mestrado-12}

\begin{enumerate}
\def\labelenumi{\arabic{enumi}.}
\item
  \href{}{Crhistian Emilio Ribeiro}. \textbf{Avaliação de redes neurais
  profundas para detecção veicular em imagens de satélite}. 2024.
  Dissertação (Engenharia de Transportes) - Universidade de São Paulo.
  Bolsa: sem bolsa\\
  DOI:
  \href{https://doi.org/10.11606/D.18.2024.tde-31102024-114415}{10.11606/D.18.2024.tde-31102024-114415}
\item
  \href{https://lattes.cnpq.br/6933325479741942}{Paola Yumi Matsumoto}.
  \textbf{Calibração de modelo Cellular Automata para simulação do
  comportamento do tráfego veicular em rodovias paulistas}. 2022.
  Dissertação (Engenharia de Transportes) - Universidade de São Paulo.
  Bolsa: CAPES\\
  DOI:
  \href{https://doi.org/10.11606/D.18.2022.tde-18102022-103557}{10.11606/D.18.2022.tde-18102022-103557}
\item
  \href{https://lattes.cnpq.br/2263620972100266}{Helena Stein Stefani}.
  \textbf{Método de previsão de fluxo de tráfego em rodovias urbanas a
  partir de dados de velocidade de fontes online}. 2021. Dissertação
  (Engenharia de Transportes) - Universidade de São Paulo. Bolsa: CNPq\\
  DOI:
  \href{https://doi.org/10.11606/D.18.2021.tde-22082022-091919}{10.11606/D.18.2021.tde-22082022-091919}
\item
  \href{https://lattes.cnpq.br/7512695029586698}{Alceu Dal Bosco
  Junior}. \textbf{Usabilidade de Pontos de Interesse e centralidades de
  rede de mapas colaborativos para análise de atração de viagens: estudo
  de caso de Curitiba}. 2020. Dissertação (Engenharia de Transportes) -
  Universidade de São Paulo. Bolsa: CAPES\\
  DOI:
  \href{https://doi.org/10.11606/D.18.2020.tde-18042022-143053}{10.11606/D.18.2020.tde-18042022-143053}
\item
  \href{https://lattes.cnpq.br/5248357327146113}{Andre Borgato Morelli}.
  \textbf{Análise exploratória de resiliência em redes viárias urbanas}.
  2019. Dissertação (Engenharia de Transportes) - Universidade de São
  Paulo. Bolsa: CNPq\\
  DOI:
  \href{https://doi.org/10.11606/D.18.2020.tde-13012020-153303}{10.11606/D.18.2020.tde-13012020-153303}
\item
  \href{https://lattes.cnpq.br/3302815314452867}{Bruna Kuramoto}.
  \textbf{Exploração de dados de mapas colaborativos em avaliações de
  morfologias urbanas brasileiras}. 2019. Dissertação (Engenharia de
  Transportes) - Universidade de São Paulo. Bolsa: CNPq\\
  DOI:
  \href{https://doi.org/10.11606/D.18.2019.tde-20082019-084513}{10.11606/D.18.2019.tde-20082019-084513}
\item
  \href{https://lattes.cnpq.br/1581979506356539}{Adriano Belletti
  Felicio}. \textbf{Identificação automática de motociclistas através de
  processamento de imagens de vídeo de tráfego}. 2019. Dissertação
  (Engenharia de Transportes) - Universidade de São Paulo. Bolsa:
  CAPES\\
  DOI:
  \href{https://doi.org/10.11606/D.18.2020.tde-12052020-170835}{10.11606/D.18.2020.tde-12052020-170835}
\item
  \href{https://lattes.cnpq.br/1828432370828925}{Gabriel Jurado Martins
  de Oliveira}. \textbf{Calibração da relação fluxo-velocidade para
  autoestradas e rodovias de pista dupla}. 2018. Dissertação (Engenharia
  de Transportes) - Universidade de São Paulo. Bolsa: CAPES\\
  DOI:
  \href{https://doi.org/10.11606/D.18.2018.tde-10092018-150848}{10.11606/D.18.2018.tde-10092018-150848}
\item
  \href{https://lattes.cnpq.br/3128200022537201}{Leandro Arab
  Marcomini}. \textbf{Identificação automática do comportamento do
  tráfego a partir de imagens de vídeo}. 2018. Dissertação (Engenharia
  de Transportes) - Universidade de São Paulo. Bolsa: CAPES\\
  DOI:
  \href{https://doi.org/10.11606/D.18.2018.tde-01102018-102649}{10.11606/D.18.2018.tde-01102018-102649}
\item
  \href{https://lattes.cnpq.br/3462901770011968}{Natália Ribeiro
  Panice}. \textbf{Método de detecção automática de eixos de caminhões
  baseado em imagens}. 2018. Dissertação (Engenharia de Transportes) -
  Universidade de São Paulo. Bolsa: CNPq\\
  DOI:
  \href{https://doi.org/10.11606/D.18.2018.tde-11122018-213600}{10.11606/D.18.2018.tde-11122018-213600}
\item
  \href{https://lattes.cnpq.br/9374994847060943}{Mariana Marçal Thebit}.
  \textbf{Reconstrução de matriz O/D sintética a partir de dados de
  tráfego disponíveis na web}. 2018. Dissertação (Engenharia de
  Transportes) - Universidade de São Paulo. Bolsa: CNPq\\
  DOI:
  \href{https://doi.org/10.11606/D.18.2018.tde-10122018-225948}{10.11606/D.18.2018.tde-10122018-225948}
\item
  \href{https://lattes.cnpq.br/4988917064689096}{Elaine Rodrigues
  Ribeiro}. \textbf{Análise exploratória de método utilizando Wavelet
  para detecção de padrões e anomalias em dados históricos do tráfego
  veicular}. 2017. Dissertação (Engenharia de Transportes) -
  Universidade de São Paulo. Bolsa: CAPES\\
  DOI:
  \href{https://doi.org/10.11606/D.18.2017.tde-07112017-212156}{10.11606/D.18.2017.tde-07112017-212156}
\end{enumerate}

\paragraph{Iniciação Científica
(16)}\label{iniciauxe7uxe3o-cientuxedfica-16}

\begin{enumerate}
\def\labelenumi{\arabic{enumi}.}
\item
  Patrick Gabriel Quintino. \textbf{Impacto da composição da frota de
  caminhões em projetos de rampas de escape de emergência}. 2025.
  Iniciação científica (Engenharia Civil) - Universidade de São Paulo.
  Bolsa: CNPq PIBIC (Programa Institucional de Bolsas de Iniciação
  Científica).
\item
  Gabriel Brunhara Alizon. \textbf{Vulnerabilidade de vias urbanas a
  eventos climáticos extremos a partir de características topográficas e
  morfológica}. 2025. Iniciação científica (Engenharia Civil) -
  Universidade de São Paulo. Bolsa: USP-EESC PUB (Programa Unificado de
  Bolsas).
\item
  Guilherme Lima Bigatao. \textbf{Análise do padrão de viagens feitas
  por bicicletas elétricas e convencionais em um sistema de
  compartilhamento de bicicletas na cidade de São Paulo}. 2024.
  Iniciação científica (Engenharia Civil) - Universidade de São Paulo.
  Bolsa: USP-EESC PUB (Programa Unificado de Bolsas).
\item
  Luan Andre Contel. \textbf{Criação de banco de dados de imagens de
  satélite para detecção de veículos}. 2024. Iniciação científica
  (Engenharia Civil) - Universidade de São Paulo. Bolsa: USP-EESC PUB
  (Programa Unificado de Bolsas).
\item
  Miguel José Sertori. \textbf{Criação de banco de dados de imagens de
  caminhões para classificação automática}. 2023. Iniciação científica
  (Engenharia Civil) - Universidade de São Paulo. Bolsa: USP-EESC PUB
  (Programa Unificado de Bolsas).
\item
  Andre de Carvalho Fiedler. \textbf{Catalogação dos processos de
  espraiamento urbano em cidades brasileiras a partir da teoria dos
  grafos}. 2022. Iniciação científica (Engenharia Civil) - Universidade
  de São Paulo. Bolsa: USP-EESC PUB (Programa Unificado de Bolsas).
\item
  Felipe Urso Parreira Pinto. \textbf{Efeito do lockdown sobre a
  pandemia de COVID-19: estudo comparativo entre Austrália e Brasil}.
  2022. Iniciação científica (Engenharia Civil) - Universidade de São
  Paulo. Bolsa: USP-EESC PUB (Programa Unificado de Bolsas).
\item
  Pedro Henrique de Lima Bertoli. \textbf{Impacto da pandemia no perfil
  de viagens dos usuários de rodovias em São Paulo}. 2022. Iniciação
  científica (Engenharia Civil) - Universidade de São Paulo. Bolsa:
  USP-EESC PUB (Programa Unificado de Bolsas).
\item
  Lucas Locatelli Helena. \textbf{Método baseado em imagem para detecção
  de eixos e classificação de caminhões}. 2021. Iniciação científica
  (Engenharia Elétrica) - Universidade de São Paulo. Bolsa: CNPq projeto
  do docente.
\item
  Danilo Bovo Carneiro. \textbf{Implementação de um modelo de simulação
  de tráfego rodoviário usando Cellular Automata}. 2020. Iniciação
  científica (Engenharia Civil) - Universidade de São Paulo. Bolsa:
  USP-EESC PUB (Programa Unificado de Bolsas).
\item
  Eraldo Dias de Castro Neto. \textbf{Análise de acessibilidade em
  ciclovias a partir do esforço total do ciclista}. 2019. Iniciação
  científica (Engenharia Civil) - Universidade de São Paulo. Bolsa: CNPq
  PIBIC (Programa Institucional de Bolsas de Iniciação Científica).
\item
  Leticia Lourenço. \textbf{Análise das alterações na legislação de
  pesos de veículos no impacto da~ vida útil do pavimento}. 2018.
  Iniciação científica (Engenharia Civil) - Universidade de São Paulo.
  Bolsa: USP-EESC PUB (Programa Unificado de Bolsas).
\item
  Luciane Sobral. \textbf{Análise do impacto da liberação do caminhão
  canavieiro de 91 toneladas na infraestrutura viária}. 2018. Iniciação
  científica (Engenharia Civil) - Universidade de São Paulo. Bolsa:
  USP-EESC PUB (Programa Unificado de Bolsas).
\item
  Paulo Cesar Rodrigues Filho. \textbf{Calibração dos parâmetros do
  comportamento dos motoristas em rodovias paulistas para o modelo de
  simulação VISSIM}. 2018. Iniciação científica (Engenharia Civil) -
  Universidade de São Paulo. Bolsa: USP-EESC PUB (Programa Unificado de
  Bolsas).
\item
  Luiza Fonseca Orlando. \textbf{Caracterização do comportamento dos
  motoristas a partir de dados compartilhados em serviços de mapeamento
  de tráfego}. 2018. Iniciação científica (Engenharia Civil) -
  Universidade de São Paulo. Bolsa: sem bolsa.
\item
  Fernando Silva Lima. \textbf{Calibração dos parâmetros fundamentais
  das curvas fluxo-velocidade em rodovias divididas do estado de São
  Paulo}. 2017. Iniciação científica (Engenharia Civil) - Universidade
  de São Paulo. Bolsa: CNPq PIBIC (Programa Institucional de Bolsas de
  Iniciação Científica).
\end{enumerate}

\paragraph{Trabalho de Conclusão de Curso
(26)}\label{trabalho-de-conclusuxe3o-de-curso-26}

\begin{enumerate}
\def\labelenumi{\arabic{enumi}.}
\item
  Guilherme Lima Bigatão. \textbf{Desenvolvimento da cadeia de
  suprimentos em obras e reformas}. 2025. Curso (Engenharia Civil) -
  Universidade de São Paulo
\item
  Guilherme Souza Araujo. \textbf{Análise da viabilidade financeira de
  uma planta de geração de energia solar na Paraíba}. 2024. Curso
  (Engenharia Civil) - Universidade de São Paulo
\item
  Miguel José Sertori. \textbf{Estudo de Capacidade Viária devido à
  Instalação de um Edifício Residencial no Município de Bebedouro --
  SP}. 2024. Curso (Engenharia Civil) - Universidade de São Paulo
\item
  Pedro Henrique de Lima Bertoli. \textbf{Estudo de viabilidade
  econômico-financeiro para a implantação de uma linha de ônibus entre
  as cidades de Lençóis Paulista/SP e São Carlos/SP}. 2024. Curso
  (Engenharia Civil) - Universidade de São Paulo
\item
  Francisco Andreson de Moura. \textbf{Aplicação de pesquisa operacional
  na operação logística de um centro de distribuição}. 2023. Curso
  (Engenharia Civil) - Universidade de São Paulo
\item
  Tainan Rodrigues Corrêa. \textbf{Proposta de programação semafórica
  isolada de tempo fixo em um cruzamento da cidade de Lençóis
  Paulista-SP}. 2023. Curso (Engenharia Civil) - Universidade de São
  Paulo
\item
  Helena Tanoue Vizioli. \textbf{Implementação de controle de velocidade
  limite variável na rodovia dos Bandeirantes}. 2022. Curso (Engenharia
  Civil) - Universidade de São Paulo
\item
  Láisla Beatriz de Carvalho Penido. \textbf{Análise da viabilidade da
  implantação de um sistema de parcel locker na cidade de Campinas/SP}.
  2021. Curso (Engenharia Civil) - Universidade de São Paulo
\item
  Gabriel Passos Bandeira. \textbf{Análise do impacto da testagem
  populacional na disseminação da COVID-19 na cidade de Ribeirão Preto
  no primeiro semestre de 2021}. 2021. Curso (Engenharia Civil) -
  Universidade de São Paulo
\item
  Breno da Cunha Costa. \textbf{Aplicação de políticas ESG em operação
  logística de empresa varejista}. 2021. Curso (Engenharia Civil) -
  Universidade de São Paulo
\item
  Eraldo Dias de Castro Neto. \textbf{Projeto de ciclovia
  energeticamente confortável ao usuário para a cidade de São Carlos}.
  2021. Curso (Engenharia Civil) - Universidade de São Paulo
\item
  Henrique Luiz Shibata Gino. \textbf{Proposta de modelo operacional
  para viagens de ônibus durante a pandemia da COVID-19}. 2021. Curso
  (Engenharia Civil) - Universidade de São Paulo
\item
  Edmar Pereira dos Santos Filho. \textbf{Redução do consumo de diesel
  em paradas na operação ferroviária por meio do desligamento automático
  de locomotivas}. 2021. Curso (Engenharia Civil) - Universidade de São
  Paulo
\item
  Jessica Tong. \textbf{Measuring the temporal and spatial impacts of
  short-term events on pedestrian flows}. 2020. Curso (Civil
  Engineering) - The University of Melbourne
\item
  Felipe Baldisseri. \textbf{Avaliação de uma nova linha metroviáriana
  região metropolitana de São Paulo baseada nos dados da Pesquisa
  Origem-Destino do Metrô SP}. 2019. Curso (Engenharia Civil) -
  Universidade de São Paulo
\item
  Eduardo Sene Eisele. \textbf{Estudo da viabilidade de implantação de
  um centro de distribuição urbana na cidade de São Carlos}. 2019. Curso
  (Engenharia Civil) - Universidade de São Paulo
\item
  Luis Gustavo Müller. \textbf{Estudo da viabilidade de implantação de
  uma rede de parcel lockers na cidade de Piracicaba}. 2019. Curso
  (Engenharia Civil) - Universidade de São Paulo
\item
  Kaique Dantas Oliveira. \textbf{Intervenção de mobilidade urbana na
  Praça Antonio Adolpho Lobbe (rotatória do Cristo) na cidade de São
  Carlos-SP}. 2019. Curso (Engenharia Civil) - Universidade de São Paulo
\item
  Francisco Mattos Fortes. \textbf{Projeto de sistema de progressão
  semafórica para as Av. Dr.~Carlos Botelho e R. 15 de Novembro}. 2019.
  Curso (Engenharia Civil) - Universidade de São Paulo
\item
  Isadora Gaidzakian Jorge. \textbf{Adequação do campus para utilização
  de veículos autônomos - estudo da ociosidade de vagas de
  estacionamento}. 2018. Curso (Engenharia Civil) - Universidade de São
  Paulo
\item
  Fernando Silva Lima. \textbf{Comparação dos níveis de serviço em
  rodovias obtidos através do método do HCM e do calibrado para rodovias
  paulistas}. 2018. Curso (Engenharia Civil) - Universidade de São Paulo
\item
  Rafael Kiyoshi Shitara. \textbf{Avaliação do plano semafórico de
  cruzamentos críticos de São Carlos}. 2016. Curso (Engenharia Civil) -
  Universidade de São Paulo
\item
  André Borgato Morelli. \textbf{Projeto de rede cicloviária para o
  município de Monte Alto, SP}. 2016. Curso (Engenharia Civil) -
  Universidade de São Paulo
\item
  Humberto Claudio Manrique. \textbf{Desenvolvimento de ferramenta de
  análise de capacidade e nível de serviço de rodovias}. 2015. Curso
  (Engenharia Civil) - Universidade de São Paulo
\item
  Guilherme Niobey Frossard. \textbf{Proposta de um~ novo terminal de
  passageiros para o aeroporto de São Carlos}. 2015. Curso (Engenharia
  Civil) - Universidade de São Paulo
\item
  Diego César Corte. \textbf{Proposta de um novo projeto geométrico para
  a interseção mais crítica da rodovia Régis Bittencourt}. 2015. Curso
  (Engenharia Civil) - Universidade de São Paulo
\end{enumerate}

\subsubsection{Em andamento}\label{em-andamento}

\paragraph{Doutorado}\label{doutorado}

\begin{enumerate}
\def\labelenumi{\arabic{enumi}.}
\item
  \href{https://lattes.cnpq.br/2619402198229137}{Thiago Vinícius Louro}.
  \textbf{Estimating the impacts of electric bicycles on accessibility
  to jobs and spatial equity}. Bolsa: CAPES. Convênio de duplo diploma:
  USP-EESC \& University Twente. Defesa: 02/04/2026
\item
  \href{}{Elaine Rodrigues Ribeiro}. \textbf{Metodologia~ para obtenção
  e análise de dados de comportamento do motociclista na corrente de
  tráfego}. Bolsa: CAPES
\item
  \href{http://lattes.cnpq.br/7503266177637246}{Pedro Henrique Caldeira
  Caliari}. \textbf{Análise da escolha modal a partir de simulação
  baseada em agentes}. Bolsa: CNPq
\end{enumerate}

\paragraph{Mestrado}\label{mestrado}

\begin{enumerate}
\def\labelenumi{\arabic{enumi}.}
\item
  \href{http://lattes.cnpq.br/7862988333397510}{Andressa Vitório Costa}.
  \textbf{Acessibilidade aos CREAS: análise da distância física e social
  dos cidadãos da região metropolitana de Belo Horizonte}. Bolsa: sem
  bolsa
\item
  \href{http://lattes.cnpq.br/6355612958124418}{Maria Eduarda Saquetto
  Michelini}. \textbf{Acessibilidade Territorial e Desigualdade: Um
  Estudo das Regionais de Curitiba}. Bolsa: sem bolsa
\item
  \href{http://lattes.cnpq.br/3660336154224847}{Rodrigo Otávio Fraga
  Peixoto de Oliveira}. \textbf{Comparação de métodos de previsão de
  alagamentos para avaliação de resiliência urbana: Estudo aplicado ao
  transporte público de Curitiba, PR}. Bolsa: sem bolsa
\item
  \href{https://lattes.cnpq.br/2158424254735233}{Murilo Colin da
  Silveira}. (tema a definir). Bolsa: CNPq
\end{enumerate}

\paragraph{Iniciação Científica}\label{iniciauxe7uxe3o-cientuxedfica}

\begin{enumerate}
\def\labelenumi{\arabic{enumi}.}
\item
  \href{https://lattes.cnpq.br/1401504528774648}{Fabricio de Matos}.
  \textbf{Análise do acesso à equipamentos públicos por ônibus em
  cidades brasileiras}. Bolsa: PIBIC
\item
  \href{\%5Bhttps://lattes.cnpq.br/1401504528774648\%5D(http://lattes.cnpq.br/4242646632967599)}{Isabelly
  Cristine Cândido Batista}. \textbf{Análise de acessibilidade por
  ônibus em Santa Rita do Sapucaí-MG}. Bolsa: sem bolsa
\end{enumerate}

\section{Produção Científica}\label{produuxe7uxe3o-cientuxedfica}

Os autores que foram orientados pelo prof. André Luiz Cunha estão
marcados em \ul{sublinhado}.

\subsection{Artigos completos publicados em periódicos
(19)}\label{artigos-completos-publicados-em-periuxf3dicos-19}

\begin{enumerate}
\def\labelenumi{\arabic{enumi}.}
\item
  \ul{LOURO, T.V.}; GRIGILON, A.B.; TIRACHINI, A.; \textbf{CUNHA, A.L.};
  GEURS, K.T. (2026) ``How Do E-Bikes Measure Up? Analyzing Speed
  Differences and Network Impacts of São Paulo's Bikesharing System''.
  Transportation.\\
  \textless DOI:
  \href{https://doi.org/10.1007/s11116-026-10724-y}{10.1007/s11116-026-10724-y}\textgreater{}
\item
  \ul{LOURO, T.V.}; GRIGILON, A.B.; \textbf{CUNHA, A.L.}; GEURS, K.T.
  (2025) ``E-bikes' impact on job accessibility and equity in São Paulo
  and Rio''. Transportation Research Part D: Transport and
  Environment.\\
  \textless DOI:
  \href{https://doi.org/10.1016/j.trd.2025.105072}{10.1016/j.trd.2025.105072}\textgreater{}
\item
  \ul{DE OLIVEIRA, G.J.M.}; LAVIERI, P.S.; \textbf{CUNHA, A.L.} (2023)
  Integrating a non-gridded space representation into a graph neural
  networks model for citywide short-term crash risk prediction. Urban
  Informatics. v.2, p.7.\\
  \textless DOI:
  \href{https://doi.org/10.1007/s44212-023-00032-6}{10.1007/s44212-023-00032-6}\textgreater{}
\item
  FLEURY, M.P.; KAMAKURA, G.K.; PITOMBO, C.S.; \textbf{CUNHA, A.L.B.N.};
  FERREIRA, F.B.; LINS DA SILVA, J. (2023) Assessing and Predicting
  Geogrid Reduction Factors after Damage Induced by Dropping Recycled
  Aggregates. Sustainability. v.15, p.9942.\\
  \textless DOI:
  \href{https://doi.org/10.3390/su15139942}{10.3390/su15139942}\textgreater{}
\item
  FLEURY, M.P.; KAMAKURA, G.K.; PITOMBO, C.S.; \textbf{CUNHA, A.L.B.N.};
  LINS DA SILVA, J. (2023) Prediction of non-woven geotextiles'
  reduction factors for damage caused by the drop of backfill materials.
  Geotextiles and Geomembranes. v.1, p.1 - 11.\\
  \textless DOI:
  \href{https://doi.org/10.1016/j.geotexmem.2023.05.004}{10.1016/j.geotexmem.2023.05.004}\textgreater{}
\item
  SILVA, F.A.E.; BESSA JUNIOR, J.E.; COSTA, A.L.; \textbf{CUNHA, A.L.};
  VELHO, D.M.C.; ANDALICIO, A. (2023) Exploratory analysis of the VISSIM
  simulation model for two-lane highways. Engenharia Civil UM (Braga),
  n.63, p.6-17. ~ \textless DOI:
  \href{https://doi.org/10.21814/ecum.4493}{10.21814/ecum.4493}\textgreater{}
\item
  SILVA, F.A.; BESSA JUNIOR, J.E.; COSTA, A.L.; \textbf{CUNHA, A.L.};
  VELHO, D.M.C. (2022) Analysis of no-passing zones to assess the level
  of service on two-lane rural highways in Brazil. Case Studies on
  Transport Policy. v.10, p.248-256.\\
  \textless DOI:
  \href{https://doi.org/10.1016/j.cstp.2021.12.006}{10.1016/j.cstp.2021.12.006}\textgreater{}
\item
  \ul{MORELLI, A. B.}; \textbf{CUNHA, A.L.} (2021) Assessing
  vulnerabilities in transport networks: a graph-theoretic approach.
  Transportes (Rio de Janeiro). v.29, p.161-172.\\
  \textless DOI:
  \href{https://doi.org/10.14295/transportes.v29i1.2250}{10.14295/transportes.v29i1.2250}\textgreater{}
\item
  SILVA, F.A.; BESSA JÚNIOR, J.E.; COSTA, A.L.; \textbf{CUNHA, A.L.};
  ANDALÍCIO, A.F.; DA COSTA VELHO, D.M.; NAZARETH, V.S. (2021)
  Evaluation of the effect of climbing lanes on segments of two-lane
  highways. Transportes (Rio de Janeiro). v.29, p.1-16.\\
  \textless DOI:
  \href{https://doi.org/10.14295/transportes.v29i2.2359}{10.14295/transportes.v29i2.2359}\textgreater{}
\item
  \ul{MORELLI, A.B.}; \textbf{CUNHA, A.L.} (2021) Measuring urban road
  network vulnerability to extreme events: An application for urban
  floods. Transportation Research Part D -- Transport and Environment.
  v.93, p.102770.\\
  \textless DOI:
  \href{https://doi.org/10.1016/j.trd.2021.102770}{10.1016/j.trd.2021.102770}\textgreater{}
\item
  MARTINS, D.O.; \ul{OLIVEIRA, G.J.M.}; MORAES, F.R.; SILVA, I.;
  \textbf{CUNHA, A.L.} (2020) Geomatics data management system. Revista
  Brasileira de Geomática. v.8, p.056-069.\\
  \textless DOI:
  \href{https://doi.org/10.3895/rbgeo.v8n1.10141}{10.3895/rbgeo.v8n1.10141}\textgreater{}
\item
  PIANUCCI, M.N.; PITOMBO, C.S.; \textbf{CUNHA, A.L.}; LIMA SEGANTINE,
  P.C. (2019) Forecasting household travel demand through a sequential
  method based on synthetic population and artificial neural networks.
  Transportes (Rio de Janeiro). v.27, p.1-23.\\
  \textless DOI:
  \href{https://doi.org/10.14295/transportes.v27i4.1409}{10.14295/transportes.v27i4.1409}\textgreater{}
\item
  OLIVEIRA, J.V.M.; LAROCCA, A.P.C.; ARAUJO NETO, J.O.; \textbf{CUNHA,
  A.L.}; SANTOS, M.C.; SCHAAL, R.E. (2019) Rigid Bridges Health Dynamic
  Monitoring Using 100 Hz GPS Single-Frequency and Accelerometers.
  Positioning. v.10, p.17-33.\\
  \textless DOI:
  \href{https://doi.org/10.4236/pos.2019.102002}{10.4236/pos.2019.102002}\textgreater{}
\item
  DE OLIVEIRA, J.V.M.; LAROCCA, A.P.C.; DE ARAÚJO NETO, J.O.; CUNHA,
  A.L.; DOS SANTOS, M.C.; SCHAAL, R.E. (2019) Vibration monitoring of a
  small concrete bridge using wavelet transforms on GPS data. Journal Of
  Civil Structural Health Monitoring. v.9, p.397-409.\\
  \textless DOI:
  \href{https://doi.org/10.1007/s13349-019-00341-y}{10.1007/s13349-019-00341-y}\textgreater{}
\item
  LINDNER, A.; PITOMBO, C.S.; \textbf{CUNHA, A.L.} (2017) Estimating
  motorized travel mode choice using classifiers: An application for
  high-dimensional multicollinear data. Travel Behaviour and Society.
  v.6, p.100-109.\\
  \textless DOI:
  \href{https://doi.org/10.1016/j.tbs.2016.08.003}{10.1016/j.tbs.2016.08.003}\textgreater{}
\item
  SOUZA, N.C.; PITOMBO, C.; \textbf{CUNHA, A.L.}; LAROCCA, A.P.C.; DE
  ALMEIDA FILHO, G.S. (2017) Model for classification of linear erosion
  processes along railways through decision tree algorithm and
  geotechnologies. Boletim de Ciências Geodésicas. v.23, p.72-86.\\
  \textless DOI:
  \href{https://doi.org/10.1590/S1982-21702017000100005}{10.1590/S1982-21702017000100005}\textgreater{}
\item
  ANDRADE, G.R.; PITOMBO, C.; \textbf{CUNHA, A.L.N.}; SETTI, J.R. (2016)
  A Model for Estimating Free-Flow Speed on Brazilian Expressways.
  Transportation Research Procedia. v.15, p.378-388.\\
  \textless DOI:
  \href{https://doi.org/10.1016/j.trpro.2016.06.032}{10.1016/j.trpro.2016.06.032}\textgreater{}
\item
  LAROCCA, A.P.C.; ARAÚJO NETO, J.O.; TRABANCO, J.L.A.; BARBOSA, A.C.B.;
  \textbf{CUNHA, A.L.B.N.}; SCHAAL, R.E. (2015) Use of 100 Hz GPS
  receivers in the detection of millimeter vertical deflections of small
  concrete bridges. Boletim de Ciências Geodésicas. v.21, p.290-307.\\
  \textless DOI:
  \href{https://doi.org/10.1590/S1982-21702015000200017}{10.1590/S1982-21702015000200017}\textgreater{}
\item
  LAROCCA, A.P.C.; ARAUJO NETO, J.O.; BARBOSA, A.C.B.; TRABANCO, J.L.A.;
  \textbf{CUNHA, A.L.B.N.} (2014) Dynamic Monitoring vertical Deflection
  of Small Concrete Bridge Using Conventional Sensors And 100 Hz GPS
  Receivers - Preliminary Results. IOSRJEN Journal of Engineering. v.04,
  p.09-20.\\
  \textless DOI:
  \href{https://doi.org/10.9790/3021-04920920}{10.9790/3021-04920920}\textgreater{}
\end{enumerate}

\subsection{Artigos completos publicados em anais de eventos
(29)}\label{artigos-completos-publicados-em-anais-de-eventos-29}

\begin{enumerate}
\def\labelenumi{\arabic{enumi}.}
\item
  \ul{MORELLI, A.B.}; \ul{ALIZON, G.B.}; \textbf{CUNHA, A.L.} (2025)
  Eficiência de alternativa em cidades brasileiras: como os padrões de
  colapso por alagamento diferem de bloqueios aleatórios. In: XXXIX
  ANPET - Congresso Nacional de Pesquisa e Ensino em Transportes, 2025,
  Goiânia. Anais do 39º ANPET.
\item
  \ul{LOURO, T.V.}; ASSIS, L.B.M.; JUNIOR, J.U.P.; \textbf{CUNHA, A.L.};
  GEURS, K.T (2025) Job Accessibility in the 15-Minute City: A
  Comparative Analysis of Walking, Cycling, and E-Bikes in Four
  Brazilian Cities. In: XXXIX ANPET - Congresso Nacional de Pesquisa e
  Ensino em Transportes, 2025, Goiânia. Anais do 39º ANPET.
\item
  ISHIHARA, B.A.; \ul{QUINTINO, P.G.}; \textbf{CUNHA, A.L.}; SETTI, J.R.
  (2025) Operação de veículos pesados em declives longos e íngremes:
  Simulação térmica de freios com base na NBR 10966-2. In: XXXIX ANPET -
  Congresso Nacional de Pesquisa e Ensino em Transportes, 2025, Goiânia.
  Anais do 39º ANPET.
\item
  RODRIGUES, L.R.; PITOMBO, C.S.; \textbf{CUNHA, A.L.}; LAROCCA, A.P.C.;
  FERRAZ, A.C.P. (2025) Identificação de hotspots de crimes e sinistros
  de trânsito no entorno de pontos de ônibus. In: XXXIX ANPET -
  Congresso Nacional de Pesquisa e Ensino em Transportes, 2025, Goiânia.
  Anais do 39º ANPET.
\item
  DAVOLI, J.P.; PITOMBO, C.S.; \textbf{CUNHA, A.L.} (2025) Aplicação de
  método de seleção de variáveis baseado em floresta aleatória para
  análise do transporte ativo pós pandemia da Covid-19. In: XXXIX ANPET
  - Congresso Nacional de Pesquisa e Ensino em Transportes, 2025,
  Goiânia. Anais do 39º ANPET.
\item
  \ul{MORELLI, A.B.}; \textbf{CUNHA, A.L.} (2024) Vulnerabilidade a
  alagamentos: como a predominância de trajetos longos diminui a
  eficiência de rotas alternativas. In: XXXVIII ANPET - Congresso
  Nacional de Pesquisa e Ensino em Transportes, 2024, Florianópolis.
  Anais do 38º ANPET.
\item
  \ul{MARCOMINI, L.A.}; \textbf{CUNHA, A.L.} (2023) Truck axle detection
  using Neural Networks: analysis of the number of images in the
  training dataset. In: ANPET - Congresso Nacional de Pesquisa e Ensino
  em Transportes, 2023, Santos. Anais do 37º ANPET.
\item
  \ul{MORELLI, A.B.}; \ul{LOURO, T.V.}; \textbf{CUNHA, A.L.} (2022)
  Proposta de indicadores de potencial cicloviário pela ótica da
  acessibilidade: identificando vias mais propícias para ciclovias e
  ciclofaixas a partir de dados amplamente disponíveis. In: XXXVI ANPET
  - Congresso Nacional de Pesquisa e Ensino em Transportes, 2022,
  Fortaleza. Anais do XXXVI ANPET.
\item
  \ul{MORELLI, A.B.}; \textbf{CUNHA, A.L.} (2021) Acessibilidade pelo
  modo a pé: impactos de características morfológicas e demográficas no
  acesso a estabelecimentos. In: XXXV ANPET - Congresso Nacional de
  Pesquisa e Ensino em Transportes, 2021. Anais do XXXV ANPET.
\item
  OLIVATTO, T.F.; PITOMBO, C.S.; \textbf{CUNHA, A.L.}; MELANDA, E.A.
  (2020) Relações entre Estado Nutricional de Pré-Escolares e
  Indicadores Socioeconômicos e de Infraestrutura Urbana: uma Aborgadem
  Baseada em Algoritmo CART. In: PLURIS - Congresso Luso-Brasileiro para
  o Planejamento Urbano, Regional, Integrado e Sustentável, 2020. Anais
  do 9º PLURIS.
\item
  \ul{BOSCO JUNIOR, A.D.}; \textbf{CUNHA, A.L.} (2020) Street and zonal
  scale relationship between network centrality and economic activities:
  case study in Curitiba, Brazil. In: PLURIS - Congresso Luso-Brasileiro
  para o Planejamento Urbano, Regional, Integrado e Sustentável, 2020,
  2020. Anais do 9º PLURIS.
\item
  \ul{VIZIOLI, H.T.}; KUŠIC, K.; IVANJKO, E.; \textbf{CUNHA, A.L.}
  (2020) A method to calibrate Variable Speed Limit Control on
  high-truck-share roads: a case study in Brazil. In: Brazilian
  Technology Symposium - BTSym'20, 2020, Campinas. Smart Innovation,
  Systems and Technologies.
\item
  CAKIJA, D.; ASSIRATI, L.; IVANJKO, E.; \textbf{CUNHA, A.L.} (2019)
  Autonomous Intersection Management: A Short Review. In: 61st
  International Symposium ELMAR-2019, Zadar. Anais do 61st ELMAR
  Symposium.
\item
  \ul{MORELLI, A.B.}; \textbf{CUNHA, A.L.} (2019) Verificação de
  vulnerabilidades em redes de transporte: uma abordagem pela teoria dos
  grafos. In: XXXIII ANPET - Congresso Nacional de Pesquisa e Ensino em
  Transportes, 2019, Balneário Camboriú. Anais do XXXIII ANPET.
\item
  \ul{KURAMOTO, B.}; \textbf{CUNHA, A.L.} (2019) Usabilidade e
  limitações de dados de mapas colaborativos em análises de
  acessibilidade. In: XXXIII ANPET - Congresso Nacional de Pesquisa e
  Ensino em Transportes, 2019, Balneário Camboriú. Anais do XXXIII
  ANPET.
\item
  SILVA, F.A.E.; BESSA JUNIOR, J.E.; COSTA, A.L.; \textbf{CUNHA, A.L.};
  ANDALICIO, A.F.; VELHO, D.M.C.; NAZARETH, V.S. (2019) Avaliação do
  efeito de faixas adicionais de subida em segmentos de rodovias de
  pista simples. In: XXXIII ANPET - Congresso Nacional de Pesquisa e
  Ensino em Transportes, 2019, Balneário Camboriú. Anais do XXXIII
  ANPET.
\item
  \ul{MARCOMINI, L.A.}; \textbf{CUNHA, A.L.} (2019) The impact of
  different video resolutions in a feature-based vehicle detection
  algorithm. In: XXXIII ANPET, 2019, Balneário Camboriú. Anais do XXXIII
  ANPET.
\item
  \ul{MORELLI, A.B.}; \textbf{CUNHA, A.L.} (2019) Uma estratégia para
  avaliação do impacto de alagamentos em sistemas viários urbanos. In:
  XXXIII ANPET - Congresso Nacional de Pesquisa e Ensino em Transportes,
  2019, Balneário Camboriú. Anais do XXXIII ANPET.
\item
  \ul{MORELLI, A.B.}; \textbf{CUNHA, A.L.} (2018) Métodos para avaliação
  de condições de tráfego a partir de dados do Google Traffic e do
  Twitter. In: XXXII ANPET, 2018, Gramado, RS. Anais do XXXII ANPET.
\item
  \ul{KURAMOTO, B.}; \textbf{CUNHA, A.L.} (2018) Proposta metodológica
  para construção e análise de uma rede real urbana. In: PLURIS -
  Congresso Luso-Brasileiro para o Planejamento Urbano, Regional,
  Integrado e Sustentável, 2018, Coimbra. Anais do 8º PLURIS.
\item
  \ul{THEBIT, M.M.}; \textbf{CUNHA, A.L.} (2017) Comparação de dados de
  tráfego disponíveis na web e obtidos por sensores fixos. In: XXXI
  ANPET, 2017, Recife. Anais do XXXI ANPET.
\item
  \ul{RIBEIRO, E.R.}; \textbf{CUNHA, A.L.} (2017) Análise exploratória
  de método para detecção de anomalias em dados de tráfego veicular
  utilizando Wavelet. In: XXXI ANPET - Congresso Nacional de Pesquisa e
  Ensino em Transportes, 2017, Recife. Anais do XXXI ANPET.
\item
  \ul{PANICE, N.R.}; \textbf{CUNHA, A.L.} (2017) Avaliação de um método
  para detecção automática de eixos de caminhões em imagens. In: XXXI
  ANPET - Congresso Nacional de Pesquisa e Ensino em Transportes, 2017,
  Recife. Anais do XXXI ANPET.
\item
  \ul{OLIVEIRA, G.J.M.}; \textbf{CUNHA, A.L.} (2017) Método de
  calibração do HCM para rodovias de pistas duplas e autoestradas a
  partir da Inferência Bayesiana. In: XXXI ANPET - Congresso Nacional de
  Pesquisa e Ensino em Transportes, 2017, Recife. Anais do XXXI ANPET.
\item
  ASSIS, R.K.M.; DURAN, J.B.C.; \textbf{CUNHA, A.L.}; PITOMBO, C.S.;
  FERNANDES JUNIOR, J.L. (2016) Aplicação de técnica de Redes Neurais
  Artificiais na criação de modelos preditivos da classificação
  funcional de pavimentos. In: PLURIS Congresso Luso-Brasileiro para o
  Planejamento Urbano, Regional, Integrado e Sustentável, 2016, Maceió.
  Anais do 7º PLURIS.
\item
  FERREIRA, F.A.; PITOMBO, C.S.; \textbf{CUNHA, A.L.} (2016) Previsão de
  escolha modal em um campus universitário utilizando a regressão
  logística binomial In: PLURIS Congresso Luso-Brasileiro para o
  Planejamento Urbano, Regional, Integrado e Sustentável, 2016, Maceió.
  Anais do 7th PLURIS.
\item
  \ul{THEBIT, M.M.}; \textbf{CUNHA, A.L.}; PITOMBO, C.S. (2016) Relação
  entre a oferta do modo ônibus e a escolha modal para acesso a
  aeroportos: uma abordagem através de data mining In: PLURIS Congresso
  Luso-Brasileiro para o Planejamento Urbano, Regional, Integrado e
  Sustentável, 2016, Maceió. Anais do 7th PLURIS.
\item
  \ul{RIBEIRO, E.R.}; \textbf{CUNHA, A.L.} (2016) Análise exploratória
  de método para definição de dia típico utilizando transformada Wavelet
  e análise de agrupamento. In: XXX ANPET - Congresso Nacional de
  Pesquisa e Ensino em Transportes, 2016, Rio de Janeiro, RJ. Anais do
  XXX ANPET.
\item
  ANDRADE, G.R.; PITOMBO, C.S.; \textbf{CUNHA, A.L.B.N.}; SETTI, J.R.;
  FERRAZ, A.C.P. (2015) Previsão da velocidade de fluxo livre em
  autoestradas e rodovias de pista dupla paulistas In: Congresso
  Brasileiro de Rodovias e Concessões -- CBR\&C, 2015, Brasília. Anais
  do 9º CBR\&C.
\item
  ROCHA, S.S.; PIANUCCI, M.N.; PITOMBO, C.S.; \textbf{CUNHA, A.L.B.N.}
  (2015) Uso de Redes Neurais para previsão de produção de viagens: uma
  análise agregada. In: XXIX ANPET - Congresso Nacional de Pesquisa e
  Ensino em Transportes, 2015, Ouro Preto, MG. Anais do XXIX ANPET.
\end{enumerate}

\subsection{Resumos publicados em anais de eventos
(6)}\label{resumos-publicados-em-anais-de-eventos-6}

\begin{enumerate}
\def\labelenumi{\arabic{enumi}.}
\item
  \ul{CALIARI, P. H. C.}; PITOMBO, CIRA; \textbf{CUNHA, A.L.} (2024)
  Avaliação da escolha modal na região metropolitana de São Paulo com
  ênfase na variabilidade espacial In: Congresso Nacional de Pesquisa e
  Ensino em Transportes (ANPET), 2024, Florianópolis. Anais do XXXVIII
  ANPET.
\item
  \ul{MORELLI, A. B.}; \textbf{CUNHA, A. L.} (2023) Incorporando
  competição na avaliação de acessibilidade dual: o método do equilíbrio
  competitivo In: Congresso de Pesquisa e Ensino em Transportes (ANPET),
  2023, Santos. Anais do 37º ANPET. \textbf{artigo vencedor do Prêmio
  ANPET}.
\item
  \ul{MORELLI, A. B.}; \textbf{CUNHA, A. L.} (2023) Detectando
  vulnerabilidade em rodovias: proposta de um indicador de redundância
  de rede para trechos da malha brasileira In: Congresso de Pesquisa e
  Ensino em Transportes (ANPET), 2023, Santos. Anais do 37º ANPET.
\item
  \ul{LOURO, T. V.}; \ul{MORELLI, A. B.}; PEDREIRA JUNIOR, J. U.;
  \textbf{CUNHA, A. L.} (2023) A metric for evaluating the potential of
  e-bikes: topographic accessibility In: Congresso de Pesquisa e Ensino
  em Transportes (ANPET), 2023, Santos. Anais do 37º ANPET.
\item
  \ul{BOSCO JUNIOR, A. D.}; \textbf{CUNHA, A. L.}. (2019) Modelos para
  atração de viagens a partir de mapas colaborativos In: XXXIII ANPET
  Congresso de Pesquisa e Ensino em Transportes, 2019, Balneário
  Camburiú.~ Anais do XXXIII ANPET.
\end{enumerate}

\subsection{Repositórios livres}\label{reposituxf3rios-livres}

O docente usou os repositórios livres para publicação \emph{preprint}
dos artigos em desenvolvimento ou com resultados preliminares. Alguns
deles foram submetidos em periódicos internacionais que aceitavam
artigos em \emph{preprint}. Um dos artigos foi publicado na
\emph{Transportation Research Part D}. Além de artigos, o docente
incentiva os orientados a disponibilizar os bancos de dados usados para
fins de comparação com outros métodos e de divulgação da pesquisa. Esses
procedimentos são muito comuns na área de Computação e adotado pelo
docente.

\subsubsection{\texorpdfstring{Artigos \emph{prepint}
(12)}{Artigos prepint (12)}}\label{artigos-prepint-12}

\begin{enumerate}
\def\labelenumi{\arabic{enumi}.}
\item
  SALINAS, K.; BARELLA, V.; \textbf{CUNHA, A.L.}; \ul{OLIVEIRA, G.J.};
  VIERA, T.; NONATO, L.G. (2025) ORDENA: ORigin-DEstiNAtion data
  exploration. arXiv: \url{https://arxiv.org/pdf/2510.18278}
\item
  \ul{MARCOMINI, L.A.}; \textbf{CUNHA, A.L.}; (2025) Using Images from a
  Video Game to Improve the Detection of Truck Axles. arXiv:
  \url{https://arxiv.org/pdf/2509.25644}
\item
  \ul{MORELLI, A. B.}; \textbf{CUNHA, A. L.} (2025) Evoluindo
  resiliência em rotas de ônibus: Proposta de um método para a
  maximização de acessibilidade em cenários de incerteza por meio de
  algoritmo genético. arXiv: \url{https://arxiv.org/pdf/2506.20375}
\item
  \ul{MORELLI, A. B.}; \textbf{CUNHA, A. L.} (2024) Incorporating
  Competition into Dual Accessibility Assessment: The Competitive
  Equilibrium Method. arXiv: \url{https://arxiv.org/pdf/2403.04879}
\item
  \ul{MORELLI, A. B.}; \textbf{CUNHA, A. L.} (2023) Resilience in
  Highways: Proposal of Roadway Redundancy Indicators and Application in
  Segments of the Brazilian Network. arXiv:
  \url{https://arxiv.org/pdf/2312.00731}
\item
  \ul{MORELLI, A. B.}; \ul{FIEDLER, A. C.}; \textbf{CUNHA, A. L.} (2023)
  Um banco de dados de empregos formais georreferenciados em cidades
  brasileiras. arXiv: \url{https://arxiv.org/pdf/2303.09602}
\item
  \ul{MARCOMINI, L.A.}; \textbf{CUNHA, A.L.}; (2023) Truck Axle
  Detection with Convolutional Neural Networks.. arXiv:
  \url{https://arxiv.org/pdf/2204.01868}
\item
  \ul{RIBEIRO, E.R.}; \textbf{CUNHA, A.L.}; (2020) Historical traffic
  flow data reconstrucion applying Wavelet Transform. arXiv:
  \url{https://arxiv.org/pdf/2006.07741}
\item
  \ul{MORELLI, A. B.}; \textbf{CUNHA, A. L.} (2020) Measuring urban road
  network resilience to extreme events: an application for urban floods.
  arXiv: \url{https://arxiv.org/pdf/1912.01739}. DOI:
  \href{https://doi.org/10.1016/j.trd.2021.102770}{10.1016/j.trd.2021.102770}
\item
  \ul{OLIVEIRA, G.J.}; \textbf{CUNHA, A. L.} (2019) A calibration method
  structured on Bayesian Inference of the HCM speed-flow relationship
  for freeways and multilane highways and a temporal analysis of traffic
  behavior. arXiv: \url{https://arxiv.org/pdf/1908.10852}
\item
  \ul{FELICIO, A.B.}; \textbf{CUNHA, A. L.} (2019) Classification of
  Motorcycles using Extracted Images of Traffic Monitoring Videos.
  arXiv: \url{https://arxiv.org/pdf/1904.00247}
\item
  \ul{MARCOMINI, L.A.}; \textbf{CUNHA, A.L.}; (2018) A Comparison
  between Background Modelling Methods for Vehicle Segmentation in
  Highway Traffic Videos. arXiv: \url{https://arxiv.org/pdf/1810.02835}
\end{enumerate}

\subsubsection{Banco de dados (6)}\label{banco-de-dados-6}

\begin{enumerate}
\def\labelenumi{\arabic{enumi}.}
\item
  \ul{MORELLI, A. B.}; \ul{ALIZON, G.B.}; \textbf{CUNHA, A. L.} (2025)
  Eficiência de Alternativa em Cidades Brasileiras: Como os Padrões de
  Colapso por Alagamento Diferem de Bloqueios Aleatórios. Zenodo:
  \url{https://zenodo.org/records/17094636}
\item
  \ul{CONTEL, L.A.}; \ul{RIBEIRO, C.E.}; \textbf{CUNHA, A.L.}; (2024)
  Criação de banco de dados de imagens de satélite para detecção de
  veículos. Zenodo: \url{https://zenodo.org/records/11214846}
\item
  \ul{RIBEIRO, C.E.}; \ul{CONTEL, L.A.}; \textbf{CUNHA, A.L.}; (2024)
  Brazilian highways dataset - satellite images. Zenodo:
  \url{https://zenodo.org/records/13992469}
\item
  \ul{MARCOMINI, L.A.}; \textbf{CUNHA, A.L.}; (2021) \# Truck Image
  Dataset. Zenodo: \url{https://zenodo.org/records/17128113}
\item
  \ul{MARCOMINI, L.A.}; \textbf{CUNHA, A.L.}; (2020) \# Truck Axle
  Detection. Zenodo: \url{https://zenodo.org/records/3788068}
\item
  \ul{MARCOMINI, L.A.}; \textbf{CUNHA, A.L.}; (2018) A Comparison
  between Background Modelling Methods for Vehicle Segmentation in
  Highway Traffic Videos. Zenodo:
  \url{https://zenodo.org/records/3612131}
\end{enumerate}

\section{Projetos de Pesquisa e
Financiamentos}\label{projetos-de-pesquisa-e-financiamentos}

\subsection{Financiamentos com avaliação por
pares}\label{financiamentos-com-avaliauxe7uxe3o-por-pares}

\begin{enumerate}
\def\labelenumi{\arabic{enumi}.}
\item
  \textbf{CCD Sustentabilidade e Inovação em Infraestrutura Rodoviária}
  -- 2025--2030\\
  Reciclagem de Pavimentos como Pilar da Descarbonização --- Centros de
  Ciência para o Desenvolvimento (CCD)\\
  \emph{Papel:} Pesquisador Associado\\
  \emph{Valor:} R\$ 2.500.000\\
  \emph{Financiador:} FAPESP (Processo
  \href{https://bv.fapesp.br/pt/auxilios/119284/ccd-sustentabilidade-e-inovacao-em-infraestrutura-rodoviaria-reciclagem-de-pavimentos-como-pilar-da-/?q=2025/07146-8}{2025/07146-8})
\item
  \textbf{Inteligência Artificial Recriando Ambientes (IARA)} --
  2023--2028\\
  Programa de Centros de Pesquisa Aplicada (CEPID)\\
  \emph{Papel:} Pesquisador Colaborador\\
  \emph{Valor:} R\$ 10.000.000\\
  \emph{Financiador:} FAPESP (Processo
  \href{https://bv.fapesp.br/pt/auxilios/111390/iara-inteligencia-artificial-recriando-ambientes/}{20/09835-1})
\item
  \textbf{Repensando a modelagem de tráfego em redes de transportes para
  a nova geração de cidades inteligentes e conectadas} -- 2023--2026\\
  \emph{Papel:} Responsável\\
  \emph{Valor:} R\$ 72.000\\
  \emph{Financiador:} CNPq --- Grupos Consolidados de Pesquisa (Processo
  409087/2023-8)
\item
  \textbf{Inteligência Artificial: desenvolvimento de ferramentas para
  mobilidade urbana} -- 2023--2026\\
  \emph{Papel:} Responsável\\
  \emph{Valor:} R\$ 40.000\\
  \emph{Financiador:} CNPq --- Bolsa de Produtividade em Pesquisa
  (Processo 311964/2022-2)
\item
  \textbf{Innovative Control Strategies for Sustainable Mobility in
  Smart Cities} -- 2021\\
  \emph{Papel:} Pesquisador Colaborador\\
  \emph{Valor:} EUR 8.000\\
  \emph{Financiador:} University of Zagreb (UNIZG), Croácia
\item
  \textbf{Uma abordagem baseada em grafos para análise espaço-temporal
  de redes viárias} -- 2020\\
  Professor Visitante Junior --- CAPES-Print/USP\\
  \emph{Papel:} Responsável\\
  \emph{Valor:} AUD 150.000\\
  \emph{Financiador:} CAPES-Print (Processo 88887.371506/2019-00)
\item
  \textbf{Método baseado em imagem para detecção de eixos e
  classificação de caminhões} -- 2018--2022\\
  \emph{Papel:} Responsável\\
  \emph{Valor:} R\$ 15.000\\
  \emph{Financiador:} CNPq --- Auxílio Universal (Processo
  436954/2018-4)
\item
  \textbf{Application of deep learning in intelligent traffic control
  system} -- 2018\\
  \emph{Papel:} Pesquisador Colaborador\\
  \emph{Valor:} EUR 8.500\\
  \emph{Financiador:} University of Zagreb (UNIZG), Croácia
\item
  \textbf{Estudos voltados à promoção da mobilidade urbana sustentável e
  segura} -- 2016--2018\\
  \emph{Papel:} Pesquisador Associado\\
  \emph{Valor:} R\$ 20.000\\
  \emph{Financiador:} CAPES/FCT (Processo 39/2014)
\end{enumerate}

\subsection{Assessoria e Consultoria}\label{assessoria-e-consultoria}

\begin{enumerate}
\def\labelenumi{\arabic{enumi}.}
\item
  \textbf{Redesenho e Validação da Rampa de Escape na BR-116 (Via
  Dutra)} -- 2025\\
  Análise do comportamento da frota de caminhões na nova pista
  descendente da Serra das Araras/RJ\\
  \emph{Papel:} Investigador Principal\\
  \emph{Valor:} R\$ 144.000\\
  \emph{Contratante:} Grupo CCR (Motiva CCR RioSP)
\item
  \textbf{Avaliação da Rampa de Escape na BR-116 (Via Dutra)} -- 2023\\
  Análise da localização de áreas de escape no trecho da Serra das
  Araras/RJ\\
  \emph{Papel:} Investigador Principal\\
  \emph{Valor:} R\$ 90.000\\
  \emph{Contratante:} Grupo CCR (CCR Nova Dutra)
\item
  \textbf{Otimização da Localização de Rampas de Escape na BR-376} --
  2019 ~ Supervisão de ensaios de homologação na área de escape para
  caminhões no km 667 da BR-376/PR\\
  \emph{Papel:} Investigador Principal\\
  \emph{Valor:} R\$ 175.631 ~ \emph{Contratante:} ARTERIS --- Autopista
  Litoral Sul (APLS)
\end{enumerate}

\subsection{Cooperações
Técnico-Científicas}\label{cooperauxe7uxf5es-tuxe9cnico-cientuxedficas}

\begin{enumerate}
\def\labelenumi{\arabic{enumi}.}
\item
  \textbf{Desenvolvimento de estudos e pesquisas em temas de
  infraestrutura de transportes e regulação} -- 2026--atual\\
  Convênio de Pesquisa entre ARTESP e USP\\
  \emph{Papel:} Pesquisador\\
  \emph{Financiador:} ARTESP (Agência de Transporte do Estado de São
  Paulo)
\item
  \textbf{Desenvolvimento de um Manual de Capacidade para rodovias
  paulistas} -- 2018--atual\\
  Termo de Cooperação Técnico-Científica ARTESP--USP/EESC para adaptação
  do HCM às condições brasileiras\\
  \emph{Papel:} Pesquisador\\
  \emph{Financiador:} ARTESP (Agência de Transporte do Estado de São
  Paulo)
\end{enumerate}

\subsection{Projetos de Extensão}\label{projetos-de-extensuxe3o}

\begin{enumerate}
\def\labelenumi{\arabic{enumi}.}
\item
  \textbf{Curso de Capacitação em Infraestrutura de Transportes e Gestão
  de Projetos} -- 2019\\
  Convênio USP--Ministério de Obras Públicas e Comunicações do Paraguai
  (MOPC).\\
  \emph{Papel:} Docente.\\
  \emph{Financiador:} MOPC
\item
  \textbf{Curso de Capacitação em Infraestrutura de Transportes} --
  2018\\
  Convênio USP--Ministério de Obras Públicas e Comunicações do Paraguai
  (MOPC).\\
  \emph{Papel:} Docente.\\
  \emph{Financiador:} MOPC
\end{enumerate}

\section{Participação em Bancas}\label{participauxe7uxe3o-em-bancas}

\subsection{Bancas de Concursos (5)}\label{bancas-de-concursos-5}

\begin{enumerate}
\def\labelenumi{\arabic{enumi}.}
\item
  Membro da banca examinadora do Concurso Público de Professor Adjunto,
  na área de Infraestrutura e Economia dos Transportes, Universidade de
  Brasília (UnB), Brasília/DF. 13 a 15 de abril de 2026
\item
  Membro da Comissão Julgadora do Concurso Público para contratação de
  Professor Doutor, na área de Engenharia de Transportes, Faculdade de
  Tecnologia Universidade Estadual de Campinas (FT-UNICAMP), Limeira/SP.
  2024
\item
  Membro da Comissão Julgadora do Concurso Público para contratação de
  Professor Doutor, na área de Engenharia de Transportes, Faculdade de
  Tecnologia Universidade Estadual de Campinas (FT-UNICAMP), Limeira/SP.
  10 e 11 de abril de 2023
\item
  Membro da banca examinadora do Concurso Público de Professor Adjunto
  A, Engenharia de Transportes, Faculdade de Engenharias e Arquitetura e
  Urbanismo e Geografia (FAENG), Campo Grande/MS. 4 a 6 de maio de 2018.
\item
  Membro da banca examinadora do Concurso Público de Professor
  Assistente A, Engenharia Civil, Infra-Estrutura de Transportes (147),
  Faculdade de Engenharias e Arquitetura e Urbanismo e Geografia
  (FAENG), Campo Grande/MS . 23 a 26 de junho de 2017.
\end{enumerate}

\subsection{Bancas de Doutorado (12)}\label{bancas-de-doutorado-12}

\begin{enumerate}
\def\labelenumi{\arabic{enumi}.}
\item
  \textbf{CUNHA, A.L.}; WIDMER, J.A.; CARVALHO, A.C.P.L.F.; CUNTO,
  F.J.C.; VASSALLO, R.F. Banca de Leandro Arab Marcomini.
  \textbf{Classificação de caminhões por eixos com aprendizagem profunda
  e um modelo multimodal de larga escala}, 2025. (Engenharia de
  Transportes) USP.
\item
  \textbf{CUNHA, A.L.}; LAVIERI, P.S.; PEREIRA, R.H.M.; CALEFFI, F.;
  MANZATO, G.G. Banca de Andre Borgato Morelli. \textbf{Análise da
  vulnerabilidade em redes urbanas brasileiras: Explorando o impacto de
  alagamentos e sua relação com a morfologia e uso do solo}, 2025.
  (Engenharia de Transportes) USP.
\item
  BERTONCINI, B.V.; \textbf{CUNHA, A.L.}; MARTE, C.L.; SILVA, G.; MEIRA,
  L.H. Banca de João Evangelista Dantas dos Santos. \textbf{Análise da
  vulnerabilidade da rede viária urbana na acessibilidade do transporte
  urbano de cargas sob a ótica da emissão de \(CO_2\)}, 2024.
  (Engenharia de Transportes) UFC.
\item
  LAVIERI, P.S.; \textbf{CUNHA, A.L.}; THOMPSON, R. Banca de Gabriel
  Jurado Martins de Oliveira. \textbf{Toward a data-driven framework to
  analyse temporal and spatial distributions of crashes}, 2024. (PhD in
  Infrastructure Engineering) The University of Melbourne.
\item
  BERTONCINI, B.V.; OLIVEIRA NETO, F.M.; BRANCO, V.T.F.C.;
  \textbf{CUNHA, A.L.}; OLIVEIRA, L.K. Banca de Demostenis Ramos
  Cassiano. \textbf{Análise integrada da influência dos aspectos dos
  sistemas de atividades e de transportes na emissão de poluentes
  decorrente da operação do transporte urbano de cargas}, 2023.
  (Engenharia de Transportes) UFC.
\item
  \textbf{CUNHA, A.L.}; CALEFFI, F.; ANZANELLO, M.J. Banca de Douglas
  Zechin. \textbf{Probabilistic Traffic Breakdown Through Bayesian
  Approximation Using Variational LSTMs}, 2023. (Engenharia de Produção)
  UFRGS.
\item
  WOLF, D.F.; VITOR, G.B.; \textbf{CUNHA, A.L.}; TOLEDO, C.F.M. Banca de
  Tiago Cesar dos Santos. \textbf{Uma Abordagem Descentralizada para
  Negociação de Tráfego de Veículos Conectados e Autônomos e Resolução
  de Conflitos}, 2023. (Ciências da Computação e Matemática
  Computacional) USP.
\item
  CYBIS, H.B.B.; WASHBURN, S.; CUNTO, F.J.C.; \textbf{CUNHA, A.L.};
  SETTI, J.R. Banca de Fernando José Piva. \textbf{Freeway level of
  service criteria based on travelers' perception}, 2022. (Engenharia de
  Transportes) USP.
\item
  SILVA, C.A.U.; PITOMBO, C.S.; \textbf{CUNHA, A.L.}; AZEVEDO FILHO,
  M.A.N.; BERTONCINI, B.V. Banca de Sameque Farias Cunha de Oliveira.
  \textbf{O uso do Twitter e Instagram para caracterização, análise de
  variabilidade e modelagem da demanda por transporte urbano em
  Fortaleza-CE}, 2020. (Engenharia de Transportes) UFC.
\item
  NOBREGA, R.A.A.; \textbf{CUNHA, A.L.}; FAVARETTO, F.; LEAL, F.;
  FERRARI, R.M.; LIMA, R.S. Banca de Harlenn dos Santos Lopes.
  \textbf{Análise do escoamento da soja brasileira através da simulação
  a eventos discretos}, 2017. (Engenharia de Produção) UNIFEI.
\item
  TRABANCO, J.L.A.; LAROCCA, A.P.C.; \textbf{CUNHA, A.L.}; COSTA, D.C.;
  FONSECA JUNIOR, E.S. Banca de João Olympio de Araújo Neto. \textbf{Uso
  do GPS no monitoramento dinâmico da infraestrutura de transportes:
  pontes rodoviárias em concreto}, 2017. (Engenharia Civil) UNICAMP.
\item
  SEGANTINE, P.C.L.; \textbf{CUNHA, A.L.}; BERTONCINI, B.V.; SIMOES,
  A.S.A.; SANCHES, S.P. Banca de Marcela Navarro Pianucci. \textbf{Uma
  proposta para obtenção da população sintética através de dados
  agregados para modelagem de geração de viagens por domicílio}, 2016.
  (Engenharia de Transportes) USP.
\end{enumerate}

\subsection{Bancas de Mestrado (33)}\label{bancas-de-mestrado-33}

\begin{enumerate}
\def\labelenumi{\arabic{enumi}.}
\item
  GIANNOTTI, M.A.; \textbf{CUNHA, A.L.}; KON, F. Banca de Flávio Soares
  de Freitas. \textbf{Acessibilidade por bicicleta às escolas públicas
  de ensino médio em São Paulo: uma abordagem qualitativa dos impactos
  da infraestrutura}, 2025. (Engenharia de Transportes) USP.
\item
  MARTE, C.L.; BESSA JÚNIOR, J.E.; \textbf{CUNHA, A.L.} Banca de Andrey
  Luis Barbosa Gomes. \textbf{Controle semafórico em tempo real para
  interseções com conflitos entre pedestres e veículos}, 2025.
  (Engenharia de Transportes) USP.
\item
  ALMEIDA FILHO, F.G.V.; \textbf{CUNHA, A.L.}; ANDRADE, M.F. Banca de
  Claudia Cosme Mascarenhas. \textbf{Desigualdades de exposição à
  poluição veicular em São Paulo: mapeamento com base em modelo de
  demanda de viagens por transporte individual motorizado}, 2025.
  (Engenharia de Transportes) USP.
\item
  GIANNOTTI, M.A.; LOUREIRO, C.F.G.; \textbf{CUNHA, A.L.} Banca de
  Thiago Silva Espindola Ferraz. \textbf{Explorando Registros de
  Detalhes de Chamadas (CDRs) para Análise da Mobilidade Urbana: Um
  Estudo Comparativo de Fortaleza, São Paulo e São José dos Campos},
  2025. (Engenharia de Transportes) USP.
\item
  ALMEIDA FILHO, F.G.V.; \textbf{CUNHA, A.L.}; CARNEIRO, M. Banca de
  João Carlos Brandão Reberte. \textbf{Análise da qualidade de produtos
  cartográficos em escalas cadastrais gerados através da técnica de
  fotogramétrica de calibração de campos mistos na câmera do VANT
  equipado com sistema Real Time Kinematic}, 2024. (Engenharia de
  Transportes) USP.
\item
  \textbf{CUNHA, A.L.}; ALMEIDA FILHO, F.G.V.; CASTRO NETO, M.M. Banca
  de Crhistian Emilio Ribeiro. \textbf{Avaliação de Redes Neurais
  Profundas para detecção veicular em imagens de satélite}, 2024.
  (Engenharia de Transportes) USP.
\item
  TACO, P.W.G.; ZANCHETTA, F.; SILVIA, R.C.; \textbf{CUNHA, A.L.} Banca
  de Kevin Masinda Mahema. \textbf{Componentes para classificação de
  corredores rodoviários inteligentes --- CRIS no Brasil}, 2023.
  (Transportes) UnB.
\item
  LIMA, R.S.; \textbf{CUNHA, A.L.}; PINHO, A.F.; FAVARETO, F. Banca de
  Lucas Gomes Pereira. \textbf{Clusterização como técnica de apoio à
  decisão para um marketplace eletrônico logístico}, 2023. (Engenharia
  de Produção) UNIFEI.
\item
  \textbf{CUNHA, A.L.}; ISLER, C.A.; BESSA JÚNIOR, J.E. Banca de Paola
  Yumi Matsumoto. \textbf{Calibração de modelo Cellular Automata para
  simulação do comportamento do tráfego veicular em rodovias paulistas},
  2022. (Engenharia de Transportes) USP.
\item
  BESSA JÚNIOR, J.E.; \textbf{CUNHA, A.L.}; GARCIA, D.S.P. Banca de Davi
  Antonio Avelar Brêtas. \textbf{Desenvolvimento de método para estimar
  o nível de serviço em interseções do tipo rotatória alongada não
  vazada localizadas em rodovias de pista simples}, 2022. (Geotecnia e
  Transportes) UFMG.
\item
  \textbf{CUNHA, A.L.}; BESSA JÚNIOR, J.E.; CASTRO NETO, M.M. Banca de
  Helena Stein Stefani. \textbf{Método de previsão de fluxo de tráfego
  em rodovias urbanas a partir de dados de velocidade de fontes online},
  2021. (Engenharia de Transportes) USP.
\item
  MELANDA, E.A.; SHIGUEMORI, E.H.; \textbf{CUNHA, A.L.}; Banca de
  Tatiane Ferreira Olivatto. \textbf{Identificação automática de rampas
  de acessibilidade apoiada por visão computacional a partir de imagens
  panorâmicas Street-Level.''}, 2021. (Engenharia Urbana) UFSCar.
\item
  SILVA JUNIOR, C.A.P.; \textbf{CUNHA, A.L.}; GUERREIRO, T.C.M.; LOPES,
  D.D.; TEIXEIRA, R.S. Banca de Thiago Vinicius Louro. \textbf{Análise
  dos fatores que influenciam a distância lateral e a velocidade de
  ultrapassagem entre veículos motorizados e ciclistas.''}, 2021.
  (Engenharia Civil) UEL.
\item
  BAZZAN, A.L.C.; SILVA, B.C.; \textbf{CUNHA, A.L.}; SILVA, R. Banca de
  Camil Samer Zahlan Redwan. \textbf{Caracterização de redes de
  transporte por meio de uma métrica de entropia}, 2020. (PPGC) UFRGS.
\item
  \textbf{CUNHA, A.L.}; MANZATO, G.G.; BERTONCINI, B.V. Banca de Alceu
  Dal Bosco Junior. \textbf{Usabilidade de pontos de interesse e
  centralidades de rede de mapas colaborativos para análise de atração
  de viagens: estudo de caso de Curitiba}, 2020. (Engenharia de
  Transportes) USP.
\item
  MARTE, C.L.; \textbf{CUNHA, A.L.}; VIANNA JUNIOR, E.O. Banca de Luca
  Di Biase. \textbf{Análise de sistemas de otimização semafórica em
  tempo real para a melhoria do desempenho da rede viária: um estudo de
  caso na cidade de São Paulo}, 2019. (Engenharia de Transportes) USP.
\item
  \textbf{CUNHA, A.L.}; LIMA, R.S.; ALMEIDA FILHO, F.G.V. Banca de André
  Borgato Morelli. \textbf{Análise exploratória de resiliência em redes
  viárias urbanas}, 2019. (Engenharia de Transportes) USP.
\item
  BESSA JUNIOR, J.E.; \textbf{CUNHA, A.L.}; ISLER, C.A. Banca de
  Frederico Amaral e Silva. \textbf{Determinação do impacto de zonas de
  ultrapassagens proibidas e de faixas adicionais em subida em segmentos
  de rodovias de pista simples}, 2019. (Geotecnia e Transportes) UFMG.
\item
  \textbf{CUNHA, A.L.}; BARROS, A.P.B.G.; LIMA, R.S. Banca de Bruna
  Kuramoto. \textbf{Exploração de dados de mapas colaborativos em
  avaliações de morfologias urbanas brasileiras}, 2019. (Engenharia de
  Transportes) USP.
\item
  \textbf{CUNHA, A.L.}; GONZAGA, A.; CUNTO, F.J.C. Banca de Adriano
  Belletti Felicio. \textbf{Identificação automática de motociclistas
  através de processamento de imagens de vídeo de tráfego}, 2019.
  (Engenharia de Transportes) USP.
\item
  MARTE, C.L.; RIBEIRO, P.C.M.; \textbf{CUNHA, A.L.} Banca de Arnaldo
  Luís Santos Pereira. \textbf{Intervenções operacionais visando a
  regularidade e a eficiência de sistemas de ônibus urbanos: resenha de
  estudos acadêmicos e simulação de aplicações com dados reais}, 2019.
  (Engenharia de Transportes) USP.
\item
  CUNTO, F.J.C.; CASTRO NETO, M.M.; \textbf{CUNHA, A.L.} Banca de Johny
  Alves Lira. \textbf{Avaliação da qualidade da extração de trajetórias
  veiculares em áreas urbanas a partir da visão computacional}, 2018.
  (Engenharia de Transportes) UFC.
\item
  \textbf{CUNHA, A.L.}; BESSA JUNIOR, J.E.; ISLER, C.A. Banca de Gabriel
  Jurado Martins de Oliveira. \textbf{Calibração da relação
  fluxo-velocidade para autoestradas e rodovias de pista dupla}, 2018.
  (Engenharia de Transportes) USP.
\item
  \textbf{CUNHA, A.L.}; CUNTO, F.J.C.; GONZAGA, A. Banca de Leandro Arab
  Marcomini. \textbf{Identificação automática do comportamento do
  tráfego a partir de imagens de vídeo}, 2018. (Engenharia de
  Transportes) USP.
\item
  \textbf{CUNHA, A.L.}; GONZAGA, A.; KRAUS JUNIOR, W. Banca de Natália
  Ribeiro Panice. \textbf{Método de detecção automática de eixos de
  caminhões baseado em imagens}, 2018. (Engenharia de Transportes) USP.
\item
  \textbf{CUNHA, A.L.}; ALMEIDA FILHO, F.G.V.; BERTONCINI, B.V. Banca de
  Mariana Marçal Thebit. \textbf{Reconstrução de matriz O/D sintética a
  partir de dados de tráfego disponíveis na web}, 2018. (Engenharia de
  Transportes) USP.
\item
  LAZZARINI, C.M.C.; \textbf{CUNHA, A.L.}; ARAUJO, F. Banca de Amilton
  da Silva Junior. \textbf{Meta-heurística baseadas na busca em
  vizinhança variável para resolução do BPP e VSBPP e aplicação no
  problema de distribuição física de café torrado e moído}, 2018.
  (Engenharia Civil) UFU.
\item
  LAZZARINI, C.M.C.; SORRATINI, J.A.; \textbf{CUNHA, A.L.} Banca de
  Gabriel José da Silva. \textbf{Microssimulação para avaliação de
  desempenho operacional da duplicação de uma rodovia: caso da BR-365},
  2017. (Engenharia Civil) UFU.
\item
  \textbf{CUNHA, A.L.}; CASTRO NETO, M.M.; ISLER, C.A. Banca de Elaine
  Rodrigues Ribeiro. \textbf{Análise exploratória de método utilizando
  Wavelet para detecção de padrões e anomalias em dados históricos do
  tráfego veicular}, 2017. (Engenharia de Transportes) USP.
\item
  CUNTO, F.J.C.; CASTRO NETO, M.M.; \textbf{CUNHA, A.L.}; CYBIS, H.B.B.
  Banca de Dênnys Araújo Santos. \textbf{Diagramas veiculares
  espaço-tempo em vias urbanas utilizando a visão computacional}, 2017.
  (Engenharia de Transportes) UFC.
\item
  BEZERRA, B.S.; \textbf{CUNHA, A.L.}; RODRIGUES, J.S. Banca de Paulo
  Roberto Alves. \textbf{Fatores na adição da tecnologia RFID no
  controle de estoque no setor sucroalcooleiro}, 2016. (Engenharia de
  Produção) UNESP.
\item
  LAZZARINI, C.M.C.; SORRATINI, J.A.; \textbf{CUNHA, A.L.} Banca de
  Jardel Inácio Moreira Vieira. \textbf{Estudo para compartilhamento de
  terminais entre um sistema tronco-alimentador de passageiros e a
  logística urbana de cargas}, 2015. (Engenharia Civil) UFU.
\item
  MULLER, C.J.; CYBIS, H.B.B.; \textbf{CUNHA, A.L.} Banca de Julia Lopes
  de Oliveira Freitas. \textbf{Modelagem de Processos para a Gestão
  Inteligente das Informações no Controle Centralizado do Tráfego},
  2014. (Engenharia de Produção) UFRGS.
\end{enumerate}

\subsection{Exames de Qualificação de Doutorado
(12)}\label{exames-de-qualificauxe7uxe3o-de-doutorado-12}

\begin{enumerate}
\def\labelenumi{\arabic{enumi}.}
\item
  PINHO, A.F.; LEAL, F.; SILVA, J.P.; \textbf{CUNHA, A.L.} Banca de José
  da Silva Ferreira Junior. \textbf{Modelo híbrido para racionalização
  do transporte público urbano por ônibus com aprendizagem de máquina},
  2025. (Engenharia de Produção) UNIFEI.
\item
  \textbf{CUNHA, A.L.}; ALMEIDA FILHO, F.G.V.; CUNTO, F.J.C. Banca de
  Elaine Rodrigues Ribeiro. \textbf{Análise do comportamento de
  motociclistas em trechos urbanos}, 2024. (Engenharia de Transportes)
  USP.
\item
  \textbf{CUNHA, A.L.}; GIANNOTTI, M.; LOUREIRO, C.F.G. Banca de Thiago
  Vinicius Louro. \textbf{Evaluation Potential Accessibility and spatial
  equity impacts of electric bicycles in São Paulo, Brazil}, 2024.
  (Engenharia de Transportes) USP.
\item
  CECY, T.M.; \textbf{CUNHA, A.L.}; SILVA, A.N.R.; SILVA, J.P.; BUGS,
  G.T. Banca de Amanda Christine Gallucci Silva. \textbf{Método para
  concepção de cenários em planos de mobilidade urbana}, 2024.
  (Sustentabilidade Ambiental Urbana) UTFPR.
\item
  CYBIS, H.B.B.; \textbf{CUNHA, A.L.}; CALEFFI, F.; ANZANELLO, M.J.
  Banca de Douglas Zechin. \textbf{Probabilistic traffic breakdown
  forecasting through Bayesian approximation using variational LSTMs},
  2023. (Engenharia de Produção) UFRGS.
\item
  CASTRO NETO, M.M.; CUNTO, F.J.C.; GOMES, J.P.P.; \textbf{CUNHA, A.L.};
  TOKUDA, E.K. Banca de Alessandro Macêdo de Araújo. \textbf{Modelagem
  do impacto das motocicletas no tráfego de interseções semaforizadas de
  Fortaleza}, 2022. (Engenharia de Transportes) UFC.
\item
  LAVIERI, P.S.; GIANNOTTI, M.A.; \textbf{CUNHA, A.L.} Banca de André
  Borgato Morelli. \textbf{Proposta de métodos e ferramentas para
  análise de vulnerabilidade em sistemas de transporte urbano pela ótica
  da acessibilidade}, 2022. (Engenharia de Transportes) USP.
\item
  \textbf{CUNHA, A.L.}; MACIEL, C.D.; KRAUS JUNIOR, W. Banca de Leandro
  Arab Marcomini. \textbf{Método de detecção e classificação de eixos de
  caminhões baseado em imagens de vídeo}, 2021. (Engenharia de
  Transportes) USP.
\item
  LAVIERI, P.S.; \textbf{CUNHA, A.L.}; THOMPSON, R.G.; NASSIR, N. Banca
  de Gabriel Jurado Martins de Oliveira. \textbf{Toward a data-driven
  framework to analyse temporal and spatial distributions of crashes},
  2020. (PhD in Infrastructure Engineering) The University of Melbourne.
\item
  WOLF, D.F.; BRANCO, K.R.L.J.C.; ROMERO, R.A.F.; \textbf{CUNHA, A.L.}
  Banca de Tiago Cesar dos Santos. \textbf{Uma abordagem descentralizada
  para negociação de tráfego e resolução de conflitos de veículos
  inteligentes}, 2018. (Ciências de Computação e Matemática
  Computacional) USP-ICMC.
\item
  FERNANDES JUNIOR, J.L.; PITOMBO, C.S.; \textbf{CUNHA, A.L.} Banca de
  Maria José Ayres Zagatto Penha. \textbf{Estudo do financiamento
  privado de rodovias: Experiências internacionais aplicadas ao Brasil},
  2016. (Engenharia de Transportes) USP.
\item
  SEGANTINE, P.C.L.; \textbf{CUNHA, A.L.}; BERTONCINI, B.V. Banca de
  Marcela Navarro Pianucci. \textbf{Modelo baseado em redes neurais para
  estimar demanda de viagens através de dados de fácil aquisição ---
  apoio ao planejamento de transportes}, 2014. (Engenharia de
  Transportes) USP.
\end{enumerate}

\subsection{Exames de Qualificação de Mestrado
(15)}\label{exames-de-qualificauxe7uxe3o-de-mestrado-15}

\begin{enumerate}
\def\labelenumi{\arabic{enumi}.}
\item
  \textbf{CUNHA, A.L.}; GIANNOTTI, M.A.; CALEFFI, F. Banca de Rodrigo
  Otávio Fraga Peixoto de Oliveira. \textbf{Comparação de métodos de
  previsão de alagamentos para avaliação de resiliência urbana: Estudo
  aplicado ao transporte público de Curitiba, PR}, 2025. (Engenharia de
  Transportes) USP-EESC.
\item
  NONATO, L\textgreater G\textgreater; \textbf{CUNHA, A.L.};
  MEDINA\textless{} J.L.P. Banca de Raissa Rosa dos Santos.
  \textbf{Explorando a relação entre a malha viária e o roubo de
  veículos em São Paulo}, 2025. (Ciências da Computação) USP-ICMC.
\item
  \textbf{CUNHA, A.L.}; FARIA, J.R.V.; ANDRADE, M.H. Banca de Maria
  Eduarda Saquetto Michelini. \textbf{Acessibilidade territorial e
  desigualdade: Um estudo das regionais de Curitiba}, 2025. (Engenharia
  de Transportes) USP-EESC.
\item
  LAROCCA, A.P.C.; \textbf{CUNHA, A.L.}; SILVA, K.C.R. Banca de Fernando
  Lima de Carvalho. \textbf{Modelo de previsão de acidentes aplicado à
  avaliação de segurança viária em rodovia brasileira}, 2025.
  (Engenharia de Transportes) USP-EESC.
\item
  CASTRO NETO, M.M.; CUNTO, F.J.C.; \textbf{CUNHA, A.L.} Banca de Daniel
  Nunes Magalhaes Arruda. \textbf{Estimação de fatores de equivalência
  veicular para motocicletas em interseções semafóricas de Fortaleza},
  2025. (Engenharia de Transportes) UFC.
\item
  \textbf{CUNHA, A.L.}; NÓBREGA, R.A.; PEDREIRA JUNIOR, J.U. Banca de
  Andressa Vitório Costa. \textbf{Acessibilidade ao CREAS: análise da
  distância física e social dos cidadãos de Belo Horizonte}, 2025.
  (Engenharia de Transportes) USP-EESC.
\item
  \textbf{CUNHA, A.L.}; CASTRO NETO, M.M.; SILVA, I. Banca de Crhistian
  Emilio Ribeiro. \textbf{Obtenção de parâmetros de tráfego a partir de
  imagens de satélite usando visão computacional}, 2022. (Engenharia de
  Transportes) USP-EESC.
\item
  BESSA JUNIOR, J.E.; \textbf{CUNHA, A.L.}; GARCIA, D.S.P. Banca de Davi
  Antônio Avelar Brêtas. \textbf{Desenvolvimento de método para cálculo
  da capacidade e do nível de serviço de interseções alternativas do
  tipo rotatória alongada não vazada em rodovias de pista simples
  brasileiras}, 2021. (Geotecnia e Transportes) UFMG.
\item
  SILVA JUNIOR, C.A.P.; \textbf{CUNHA, A.L.}; GUERREIRO, T.C.AM. Banca
  de Thiago Vinicius Louro. \textbf{Análise dos fatores que influenciam
  na ocorrência de ultrapassagens perigosas entre veículos motorizados e
  ciclistas}, 2020. (Engenharia Civil) UEL.
\item
  MARTE, C.L.; \textbf{CUNHA, A.L.}; CASTRO NETO, M.M. Banca de Olimpio
  Mendes de Barros. \textbf{Caracterização das condições de tráfego em
  tempo próximo ao real para uso em sistemas de previsão de tráfego em
  cidades de grande porte}, 2019. (Engenharia de Transportes) USP-EP.
\item
  \textbf{CUNHA, A.L.}; PITOMBO, C.S.; BERTONCINI; B.V. Banca de Alceu
  Dal Bosco Junior. \textbf{Proposta de modelo de atração de viagens a
  partir de mapas colaborativos}, 2019. (Engenharia de Transportes)
  USP-EESC.
\item
  BESSA JUNIOR, J.E.; \textbf{CUNHA, A.L.}; ISLER, C.A. Banca de
  Frederico Amaral e Silva. \textbf{Determinação do impacto de zonas de
  ultrapassagem proibidas e de faixas adicionais de subida em segmentos
  de rodovias de pista simples}, 2018. (Geotecnia e Transportes) UFMG.
\item
  \textbf{CUNHA, A.L.}; SILVA, A.N.R.; MINGHIM, R. Banca de Bruna
  Kuramoto. \textbf{Análise de aspectos da mobilidade urbana de uma rede
  real a partir de propriedades topográficas ponderadas}, 2017.
  (Engenharia de Transportes) USP-EESC.
\item
  \textbf{CUNHA, A.L.}; GONZAGA, A.; ISLER, C.A. Banca de Adriano
  Belletti Felicio. \textbf{Análise do comportamento de motociclistas
  através de sistema de processamento de imagens de vídeo}, 2017.
  (Engenharia de Transportes) USP-EESC.
\item
  FERNANDES JUNIOR, J.L.; \textbf{CUNHA, A.L.}; FABBRI, G.T.P. Banca de
  Danilo Rinaldi Bisconsini. \textbf{O uso de smartphones para a
  avaliação da irregularidade longitudinal de pavimentos}, 2015.
  (Engenharia de Transportes) USP-EESC.
\end{enumerate}

\subsection{Bancas de Trabalho de Conclusão de Curso
(Graduação)}\label{bancas-de-trabalho-de-conclusuxe3o-de-curso-graduauxe7uxe3o}

Até o presente momento foram 85 bancas de Trabalho de conclusão de curso
de graduação em Engenharia Civil, na USP-EESC. A Tabela @tab:tcc traz o
resumo quantitativo por ano.

\begin{longtable}[]{@{}ccl@{}}
\caption{Participação em bancas de Trabalho de Conclusão de Curso na
graduação em Engenharia Civil (USP-EESC), por
ano.}\label{tab:tcc}\tabularnewline
\toprule\noalign{}
Ano & Nº de bancas & Instituições \\
\midrule\noalign{}
\endfirsthead
\toprule\noalign{}
Ano & Nº de bancas & Instituições \\
\midrule\noalign{}
\endhead
\bottomrule\noalign{}
\endlastfoot
2014 & 5 & USP-EESC \\
2015 & 7 & USP-EESC \\
2016 & 13 & USP-EESC \\
2017 & 10 & USP-EESC \\
2018 & 7 & USP-EESC \\
2019 & 12 & USP-EESC \\
2020 & 5 & USP-EESC \\
2021 & 7 & USP-EESC \\
2022 & 3 & USP-EESC \\
2023 & 3 & USP-EESC \\
2024 & 5 & USP-EESC \\
2025 & 8 & USP-EESC \\
\textbf{Total} & \textbf{85} & \\
\end{longtable}

\end{document}
